\documentclass[11pt]{article}
\usepackage{amsmath,amssymb,amsthm}
\setlength{\textwidth}{6.5in}
\setlength{\textheight}{9in}
\setlength{\oddsidemargin}{0in}
\setlength{\evensidemargin}{0in}
\setlength{\topmargin}{-0.5in}
\newcommand{\url}[1]{\texttt{\detokenize{#1}}}
\newtheorem{theorem}{Theorem}[section]
\newtheorem{lemma}[theorem]{Lemma}
\newtheorem{proposition}[theorem]{Proposition}
\newtheorem{corollary}[theorem]{Corollary}
\newcommand{\F}{\mathbf F}
\newcommand{\T}{\mathbf T}
\newcommand{\Z}{\mathbf Z}
\newcommand{\Real}{\operatorname{Re}}
\newcommand{\cl}{\operatorname{cl}}
\title{Prime Profiles and Oscillation over Finite Fields\\
\large Optimal deletion, Frobenius waves, and ascent geometry}
\author{Research working draft}
\date{29 October 2026}
\begin{document}
\maketitle
\begin{abstract}
We study the minimum deletion of closed points needed to keep a cumulative
degree-weighted prime count below its geometric main term. For a smooth
projective geometrically connected curve over a fixed finite field, the
greedy deletion process has a full asymptotic description by a damped
envelope of its Frobenius phases. Either the greedy deletion set is finite,
or its count through degree \(n\) is of order \(q^{n/2}/n\). In the second
case its count is asymptotically optimal among all admissible retained sets.
Prime squares determine a parity bias; prime cubes settle the boundary
where the leading phase function has maximum zero. For elliptic curves
we classify every case, including the fields with two and three elements.
We also transfer the exact factor-density ascent identity to
\(\F_q[T]\), classify all unimodal ranks, and determine the eventual
number of modes as \(q\) grows at fixed rank. Quadratic polynomial
sampling supplies character-sensitive profiles with an explicit
square-root-accurate reference clock. At fixed odd \(q\) and rank
\(k\geq2\), these quadratic profiles realize every finite ascent word;
their maximal peak count among \(\deg D\leq H\) is of sharp order
\(\log H\), although their Galois group is always the same.
\end{abstract}

\paragraph{Status.}
This is a new finite-field manuscript containing ordinary mathematical
arguments, not yet Lean formalizations. Existing maintained proofs supply
the abstract greedy record mechanism and damped-wave method; their
integer-prime consumers do not automatically constitute finite-field
adapters. The curve zeta formula and Hasse--Weil theorem are classical
inputs. Contributions under investigation are the optimal deletion
problem, its full past-wave law, its elliptic classification, and the
complete factor-profile unimodality and large-field mode classifications,
finite-word realization, and logarithmic quadratic-profile complexity.
The character-clock estimate below is an application of classical
finite-field RH, not an independent novelty claim.
Priority and publishability remain under review. The finite computations
below are diagnostics rather than universal proofs. No new number-field
zero-free exponent is claimed.

\section{The retained set and its optimality problem}
Fix a prime power \(q\geq2\) and a smooth projective geometrically connected
curve \(C/\F_q\) of genus \(g\). Primes mean closed points. Let \(B_d\)
count its degree-\(d\) closed points, and put
\[
 N_m=\#C(\F_{q^m}),\quad
 \Theta(n)=\sum_{d=1}^n dB_d,\quad
 T(n)=\sum_{j=1}^nq^j=\frac{q^{n+1}-q}{q-1}.
\]
All cumulative quantities are zero at \(n=0\). A prime of degree \(d\)
has norm \(q^d\), so the weight \(d\) is log norm divided by \(\log q\).
A retained set \(S\) is admissible when
\[
 \Theta_S(n)=\sum_{\substack{P\in S\\\deg P\leq n}}\deg P\leq T(n)
 \quad\text{for every }n\geq1.
\]
Its deleted mass and count are
\[
 D_S(n)=\sum_{\substack{P\notin S\\\deg P\leq n}}\deg P,\qquad
 K_S(n)=\#\{P\notin S:\deg P\leq n\}.
\]
Process points in increasing degree, retaining a degree-\(d\) point
exactly when its addition preserves the budget \(T(d)\). Any order
within one degree gives the same number retained there. Write \(D,K\)
for the greedy deleted mass and count.

\begin{lemma}[Exact record sandwich]\label{lem:record}
For \(E(n)=\Theta(n)-T(n)\) and
\(R(n)=\max(0,E(1),\ldots,E(n))\), one has
\[
 R(n)\leq D(n)\leq R(n)+n.                         \tag{1}
\]
Every admissible \(S\) satisfies \(D_S(n)\geq R(n)\).
\end{lemma}
\begin{proof}
Admissibility gives \(D_S(j)\geq E(j)\), and deleted mass is
nondecreasing. This proves the lower bounds. If the greedy rule has
made no rejection, \(D=R=0\). Otherwise its last rejection through
degree \(n\) is at a point of degree \(d\leq n\). If \(m\) is deleted
mass before rejection and \(e\) the full mass up to that point minus
\(T(d)\), rejection says \(m<e\). Afterwards the mass is \(m+d<e+d\).
The full excess increases within a degree, so \(e\leq E(d)\leq R(n)\).
There is no subsequent increase of deleted mass, proving the upper bound.
\end{proof}
The generic mechanism is the one already proved by
\texttt{clockDeletedMass\_record\_bounds} in the maintained
\texttt{GreedyRegulator.lean}; this application uses degree weights.

\section{All prime powers and the Frobenius wave}
The classical zeta formula and finite-field RH give, with multiplicities,
\[
 N_m=q^m+1-\sum_{i=1}^{2g}\alpha_i^m,\qquad
 |\alpha_i|=\sqrt q.                              \tag{2}\label{eq:weil}
\]
The roots are conjugate-paired; for genus zero the sum is empty.
See Milne~\cite{Milne}. Write \(A(m)=\sum_{j=1}^mN_j\).

\begin{lemma}[Exact inversion]\label{lem:mobius}
\[
 N_m=\sum_{d\mid m}dB_d,\qquad
 \Theta(n)=\sum_{k=1}^n\mu(k)A(\lfloor n/k\rfloor). \tag{3}\label{eq:mobius}
\]
\end{lemma}
\begin{proof}
A degree-\(d\) closed point contributes \(d\) rational points over
\(\F_{q^m}\) if \(d\mid m\), and none otherwise. In the second expression
the coefficient of \(dB_d\) is
\(\sum_{v\leq n/d}\sum_{k\mid v}\mu(k)=1\).
\end{proof}

Put \(r=\sqrt q,\ z_i=\alpha_i/r,\ c_i=\alpha_i/(\alpha_i-1)\).
All denominators are nonzero because \(r>1\). For \(n\in\Z\), put
\[
 b_n=\begin{cases}q/(q-1)&n\text{ even},\\
 r/(q-1)&n\text{ odd},\end{cases}
 \qquad a_n=-\Real\sum_i c_i z_i^n-b_n.             \tag{4}\label{eq:wave}
\]

\begin{proposition}[The boundary retains prime cubes]\label{prop:expansion}
For fixed \(q,g\),
\[
 E(n)=q^{n/2}a_n-T(\lfloor n/3\rfloor)
       +O_{q,g}(n^2q^{n/4}),                      \tag{5}\label{eq:cubes}
\]
and consequently \(E(n)=q^{n/2}a_n+O_{q,g}(q^{n/3})\).
\end{proposition}
\begin{proof}
Geometric summation in (2) gives the exact identity
\[
 A(m)=T(m)+m-\sum_i c_i(\alpha_i^m-1).
\]
In (3) keep \(k=1,2,3\) with signs \(+,-,-\).
The fourth term vanishes since \(\mu(4)=0\). The entire remaining
sum is \(\sum_{k=5}^n\mu(k)A(\lfloor n/k\rfloor)\). Since
\(A(m)=O_{q,g}(q^m)\), that remainder is \(O_{q,g}(nq^{n/5})\).
The nonconstant Frobenius terms from \(k=2,3\) are
\(O_{q,g}(q^{n/4}+q^{n/6})\); linear and constant contributions are
\(O_{q,g}(n)\). Finally
\(T(\lfloor n/2\rfloor)=b_nq^{n/2}-q/(q-1)\).
These identities prove (5), retaining the square bias, every leading
zero phase, and the negative cube term. The last assertion follows
because \(n^2q^{n/4}=O_q(q^{n/3})\).
\end{proof}

\section{A full past-wave law}
Let
\[
 H=\cl\{(z_1^n,\ldots,z_{2g}^n,(-1)^n):n\in\Z\}
       \subseteq\T^{2g}\times\{1,-1\}.
\]
This is a compact group. Let \(h_n=(z_1^n,\ldots,z_{2g}^n,(-1)^n)\),
and let \(\tau\) be translation by \(h_1\). Define on \(H\)
\[
 b(\eta)=\frac{q+r}{2(q-1)}
              +\eta\,\frac{q-r}{2(q-1)},\quad
 a(w,\eta)=-\Real\sum_i c_iw_i-b(\eta),
\]
so \(a(h_n)=a_n\), and define
\[
 F(h)=\sup_{\ell\geq0}r^{-\ell}
             \max(0,a(\tau^{-\ell}h)).            \tag{6}\label{eq:envelope}
\]
This formula keeps every lag of the past wave.

\begin{theorem}[Envelope and count law]\label{thm:wave}
The function \(F\) is continuous, and
\[
 \frac{D(n)}{q^{n/2}}=F(h_n)+O_{q,g}(q^{-n/6}),\qquad
 \frac{nK(n)}{q^{n/2}}=F(h_n)+O_{q,g}(1/n).       \tag{7}\label{eq:law}
\]
No linear independence hypothesis is needed.
\end{theorem}
\begin{proof}
Truncating (6) to \(0\leq\ell\leq L\) changes it by at most
\(\|a\|_\infty r^{-L-1}\), uniformly. A uniform limit of these
continuous finite maxima is continuous.
Proposition~\ref{prop:expansion} gives
\[
 \left|\frac{R(n)}{r^n}
  -\max_{1\leq j\leq n}r^{j-n}(a_j)_+\right|
       =O_{q,g}(q^{-n/6}).
\]
The additional terms \(\ell\geq n\) in (6) contribute
\(O_{q,g}(r^{-n})\). Lemma~\ref{lem:record} changes normalized mass
by at most \(n/r^n=O_q(q^{-n/6})\), proving the first formula.
Finite summation by parts gives
\[
 K(n)=\frac{D(n)}n+
        \sum_{j=1}^{n-1}\frac{D(j)}{j(j+1)}.       \tag{8}\label{eq:Abel}
\]
Since \(D(j)=O_{q,g}(r^j)\), splitting at \(n/2\) bounds the sum
by \(O_{q,g}(r^n/n^2)\). This proves the count formula.
\end{proof}

\begin{theorem}[Finite or square-root, with no intermediate case]\label{thm:dichotomy}
If \(\max_H a\leq0\), the greedy deletion set is finite.
If \(\max_H a>0\), there exist \(c,C>0,n_0\), depending on the curve
and field, such that for every \(n\geq n_0\),
\[
 c\,q^{n/2}/n\leq K(n)\leq C\,q^{n/2}/n.          \tag{9}
\]
Every admissible retained set satisfies the same lower bound,
possibly with a different positive \(c\).
\end{theorem}
\begin{proof}
When \(\max_Ha\leq0\), (5) gives
\(E(n)\leq-T(\lfloor n/3\rfloor)+O_{q,g}(n^2q^{n/4})<0\)
eventually, because \(q^{n/3}\) dominates \(n^2q^{n/4}\).
All within-degree excesses are at most the negative endpoint excess,
so no later rejection is possible. This includes the boundary
\(\max_Ha=0\).

Otherwise choose \(\delta>0\) so \(U=\{h:a(h)>\delta\}\) is nonempty.
Translation by \(\tau\) is minimal on \(H\): its cyclic subgroup
is dense by the definition of \(H\). Every backward orbit meets \(U\),
so the open sets \(\tau^\ell U\), \(\ell\geq0\), cover \(H\).
Compactness yields a finite subcover with indices at most \(L\).
Thus \(F(h)\geq r^{-L}\delta>0\) for all \(h\), and (7) proves (9).
For arbitrary admissible \(S\),
\(D_S(n)\geq R(n)\gg q^{n/2}\) by the same record formula.
Each deletion through degree \(n\) has weight at most \(n\), hence
\(K_S(n)\geq D_S(n)/n\).
\end{proof}

\begin{corollary}[Asymptotically optimal number of deletions]\label{cor:optimal}
In the positive case let \(M(n)=\min_S K_S(n)\), with the minimum over
all admissible retained sets. Then
\[
 0\leq K(n)-M(n)=O_{q,g}(q^{n/2}/n^2),\qquad
 K(n)/M(n)=1+O_{q,g}(1/n).
\]
\end{corollary}
\begin{proof}
The minimum exists: the family contains the greedy set and counts
are nonnegative integers. Every \(S\) has \(K_S(n)\geq R(n)/n\).
Equations (1) and (8) give
\(K(n)=R(n)/n+O_{q,g}(q^{n/2}/n^2)\).
The positive case has \(M(n)\gg q^{n/2}/n\), allowing division.
\end{proof}

\begin{corollary}[Distribution with dependent phases]
For every continuous \(G:\mathbb R\to\mathbb R\),
\[
 \lim_{N\to\infty}\frac1N\sum_{n=1}^N
 G(nK(n)/q^{n/2})=\int_HG(F(h))\,dm_H(h),
\]
where \(m_H\) is normalized Haar measure.
\end{corollary}
\begin{proof}
The averages of each nontrivial character along powers of \(h_1\)
are geometric sums tending to zero. Characters approximate continuous
functions on the compact abelian group \(H\), proving equidistribution.
Apply this to \(G\circ F\), then use (7) and uniform continuity on
a compact interval containing all normalized counts.
\end{proof}

\section{Complete classification for elliptic curves}
Let \(t=q+1-\#E(\F_q)\) and let \(\alpha,\bar\alpha\) be the roots of
\(X^2-tX+q\). Put \(z=\alpha/r,\ c=\alpha/(\alpha-1)\). Then
\[
 a_n=-2\Real(cz^n)-b_n,\qquad
 |c|^2=\frac q{q+1-t},\quad
 \Real c=\frac{q-t/2}{q+1-t}.                     \tag{10}\label{eq:elliptic}
\]

\begin{theorem}[Every field, including the small exceptions]\label{thm:elliptic}
There exists an admissible retained set with finite complement
if and only if
\[
 \boxed{\,t=2\sqrt q\quad\text{or}\quad
 (q,t)\in\{(2,-2),(2,-1),(2,0),(3,0)\}.\,}         \tag{11}
\]
The first case requires \(q\) to be a square. In all other cases
the optimal deletion count is of order \(q^{n/2}/n\), with
the asymptotic optimality in Corollary~\ref{cor:optimal}.
\end{theorem}
\begin{proof}
By density and continuity, \(\max_Ha\leq0\) means exactly
\(a_n\leq0\) for every \(n\).
Suppose first \(q\geq4\), so \(r\geq2\).
If \(z=1\), then \(t=2r\) and \(c=r/(r-1)>0\), making every
\(a_n\) negative. If \(z=-1\), an odd \(n\) gives
\[
 a_n=\frac{2r}{r+1}-\frac r{q-1}
       =\frac{r(2r-3)}{(r-1)(r+1)}>0.
\]
If \(z\) has order four, take \(z=i\); direct calculation yields
\[
 a_2=\frac{q(q-3)}{(q-1)(q+1)}>0.
\]

For every remaining nonreal \(z\), odd powers either are dense on the
circle, run over all roots of an odd order \(h\geq3\), or run over a
rotated regular \(h/2\)-gon for even \(h\geq6\). The nearest point to
a prescribed angle is within \(\pi/3\). Thus some odd \(n\), or an
arbitrarily close approximation in the dense case, has
\(-2\Real(cz^n)\geq |c|\).
For nonreal \(\alpha\), \(|\alpha-1|<r+1\) strictly, so
\[
 |c|>\frac r{r+1}\geq\frac r{q-1}=b_n
 \quad(n\text{ odd}).
\]
This produces a positive phase, finishing all \(q\geq4\).

For \(q=2\), Hasse's bound leaves traces \(-2,-1,0,1,2\).
At \(t=-2\), \(2|c|=2\sqrt{2/5}<\sqrt2=\min b_n\);
at \(t=-1\), \(2|c|=\sqrt2=\min b_n\). These give nonpositive
waves. At \(t=0\), \(z=i\); the four classes of \(n\bmod4\) give
negative values, since classes zero and one are negative and
\[
 a_2=\frac{q(q-3)}{(q-1)(q+1)},\qquad
 a_3=\frac{\sqrt q(q-3)}{(q-1)(q+1)}.
\]
For \((q,t)=(2,1)\), \(z\) has infinite order. Otherwise
\(z+z^{-1}\) would be an algebraic integer, making its square
\(t^2/q=1/2\) a rational algebraic integer, a contradiction.
Odd powers are dense, and \(2|c|=2>\sqrt2\) gives a positive phase.
For \((2,2)\), \(\alpha=1+i,\ c=1-i\); explicitly \(a_5=\sqrt2>0\).

For \(q=3\), the possible integer traces are \(-3,-2,-1,0,1,2,3\).
When \(t=0\) the preceding four-class calculation gives \(a_n\leq0\),
with equality in classes two and three. At \(t=-2,-1,1,2\), the
rational number \(t^2/3\) is not an integer, proving infinite order
of \(z\) as above; moreover
\(2|c|\geq\sqrt2>\sqrt3/2=\min b_n\).
The odd orbit supplies a positive phase. For the two remaining
traces, direct evaluation gives
\[
 a_5=\sqrt3/2>0\ (t=3),\qquad
 a_{11}=3\sqrt3/14>0\ (t=-3).
\]
The dichotomy now proves all assertions, including the nonpositive
boundary cases.
\end{proof}

\section{Exact ascent transfer for irreducible polynomials}

This section concerns \(\F_q[T]\), with its affine irreducible polynomials,
rather than all closed points of a projective curve. Enumerate the monic
irreducibles \(P_1,P_2,\ldots\) in nondecreasing degree, with arbitrary
tie breaking, and write \(Q_i=q^{\deg P_i}\), \(w_i=(Q_i-1)^{-1}\).
For \(k\geq1\), let \(d_{q,k}(i)\) be the limiting probability that
\(P_i\) is the \(k\)-th distinct irreducible factor of a uniformly
sampled monic polynomial of degree \(M\), as \(M\to\infty\).
A polynomial with fewer than \(k\) distinct factors contributes to none
of these events. Multiplicity does not change the position of a factor.

\begin{lemma}[Actual sampling and the exact ascent identity]\label{lem:polyAscent}
With \(e_j\) the elementary symmetric polynomial,
\[
 d_{q,k}(i)=\frac1{Q_i}\prod_{j<i}(1-Q_j^{-1})
                  e_{k-1}(w_1,\ldots,w_{i-1}).       \tag{12}
\]
This probability is already exact whenever
\(M\geq\sum_{j=1}^i\deg P_j\). The numerical sequence does not depend
on the tie breaking. For \(k\geq2,\ i\geq k\), put
\[
 R_{k-1}(i)=
 \frac{e_{k-2}(w_1,\ldots,w_{i-1})}
      {e_{k-1}(w_1,\ldots,w_{i-1})}.
\]
All denominators are positive, and
\[
 \frac{d_{q,k}(i+1)}{d_{q,k}(i)}
   =\frac{Q_i-1+R_{k-1}(i)}{Q_{i+1}}.
\]
Consequently
\[
 d_{q,k}(i+1)>d_{q,k}(i)
 \quad\Longleftrightarrow\quad
 R_{k-1}(i)>Q_{i+1}-Q_i+1.                          \tag{13}\label{eq:polyStep}
\]
\end{lemma}
\begin{proof}
CRT gives independence of the divisibility indicators for the first
\(i\) irreducibles, with probabilities \(1/Q_j\). Every residue class
modulo \(\prod_{j\leq i}P_j\) occurs equally often among monic
degree-\(M\) polynomials when \(M\) is at least its degree. Extracting
the probability of exactly \(k-1\) earlier divisors proves (12).
Only the ordered norms enter the formula, proving tie invariance.
Finally \(e_r(w_1,\ldots,w_i)=e_r(w_1,\ldots,w_{i-1})
+w_i e_{r-1}(w_1,\ldots,w_{i-1})\); division of (12) at successive
indices proves (13).
\end{proof}

Equation (13) is the literal counterpart of the maintained rational-prime
ratio-versus-gap criterion. Inside a degree block its threshold is one.
At a jump from degree \(n\) to \(n+1\), the threshold is
\((q-1)q^n+1\). Equal norm does not erase the ordinary density ascent;
it changes the threshold and creates a block structure.

\begin{lemma}[Two coefficient bounds]\label{lem:ratioBounds}
For \(m\geq r\geq1\) positive weights, total weight \(\mu\), and
minimum weight \(w_{\min}\),
\[
 \frac r\mu\leq\frac{e_{r-1}}{e_r}
      \leq\frac r{(m-r+1)w_{\min}}.                 \tag{14}
\]
Writing \(I_q(d)\) for the number of monic irreducibles of degree \(d\),
and \(M_n=\sum_{d\leq n}I_q(d)\), one has
\[
 \mu_n:=\sum_{d\leq n}\frac{I_q(d)}{q^d-1}
      \leq H_n+\frac1{q-1}\leq H_n+1,\qquad
 M_n\geq2n\quad(n\geq4).                           \tag{15}
\]
Also \(I_q(d)\geq2\) for \(d\geq3\).
\end{lemma}
\begin{proof}
The exact identity
\[
 r e_r=\sum_{|J|=r-1}w_J\sum_{j\notin J}w_j
\]
bounds the inner sum between \((m-r+1)w_{\min}\) and \(\mu\),
giving (14). The identity \(q^d=\sum_{e\mid d}eI_q(e)\) gives
\(I_q(1)=q\), and \(I_q(d)\leq(q^d-q)/d\) for \(d\geq2\).
Hence \(I_q(d)/(q^d-1)<1/d\) for \(d\geq2\), proving the
weight bound.

For completeness, subtracting all possible proper subfield contributions
gives
\[
 d I_q(d)\geq q^d-\sum_{e=1}^{\lfloor d/2\rfloor}q^e
            \geq q^d-2q^{d/2}.
\]
For \(d\geq4\), \(2^d-2^{d/2+1}\geq2d\): check \(d=4,5\)
and then increase \(d\) by two. The expression \(x^2-2x\) increases
for \(x\geq1\), so \(I_q(d)\geq2\) for \(d\geq4\).
The formula \(I_q(3)=(q^3-q)/3\) gives the remaining degree.
Finally
\[
 I_q(1)\geq2,\quad I_q(2)=(q^2-q)/2\geq1,
 I_q(3)\geq2,\ I_q(4)=(q^4-q^2)/4\geq3
\]
give \(M_4\geq8\); each further degree adds at least two.
\end{proof}

\begin{theorem}[Complete finite-field unimodality classification]\label{thm:polyUnimodal}
For every prime power \(q\geq2\) and every integer \(k\geq1\),
the sequence \(i\mapsto d_{q,k}(i)\) is unimodal if and only if
\[
                         k\leq2.                 \tag{16}
\]
For \(k=1\) it is strictly decreasing. For \(k=2\) it increases
strictly to the last linear irreducible, then decreases strictly.
For every \(k\geq3\) there is an explicit strict drop followed by
a strict rise.
\end{theorem}
\begin{proof}
Unimodal means nondecreasing up to one index and nonincreasing
afterwards; a strict drop followed later by a strict rise rules this out.
The first \(k-1\) entries are zero, and all later entries are positive.

For \(k=1\), the step ratio is \((Q_i-1)/Q_{i+1}<1\).
For \(k=2\), the \(j\)-th linear irreducible has, for \(2\leq j\leq q\),
\[
 d_{q,2}(j)=\frac{j-1}{q^2}(1-1/q)^{j-2}.
\]
Successive linear entries increase exactly when \(j<q\).
At the last linear irreducible the ratio to the first quadratic
is \(1/q<1\). Every subsequent prefix has total weight greater than
one, so \(R_1=1/\mu<1\), and (13) forces strict decrease forever.

Take \(k\geq5\), put \(r=k-1\), and choose degree \(n=r\geq4\).
At the last degree-\(n\) irreducible, the preceding prefix has
\(m=M_n-1\) weights and minimum weight at least \((q^n-1)^{-1}\).
Since \(M_n\geq2r\), (14) gives \(R_r\leq q^n-1\).
The next norm is \(q^{n+1}\), so (13) gives a strict drop.
Before the first degree-\(n+1\) irreducible, (15) and
\(H_n+1<n\) for \(n\geq3\) give
\(R_r\geq r/\mu_n>1\).
There are at least two irreducibles of degree \(n+1\), so the
next step is a strict rise. The harmonic inequality follows from
\(H_3+1=17/6<3\) and induction.

For \(k=4,\ q\geq3\), the same proof uses \(n=r=3\):
\(M_3\geq3+3+8=14\geq6\), and \(\mu_3<3\).
For \(k=4,\ q=2\), use the degree-three boundary. Before its
last point the weights are \(1,1,1/3,1/7\), giving
\(R_3=e_2/e_3=7/2<9=16-8+1\), a strict drop.
After completing that block, \(\mu_3=55/21<3\); (14) gives
a rise between the first two degree-four points.

For \(k=3,\ q\geq3\), the ratio before the last linear
irreducible is \(2(q-1)/(q-2)\leq4\), below the jump threshold
\(q^2-q+1\geq7\). After all \(q\) linear points,
\(R_2=e_1/e_2=2>1\), giving a rise between the first two
quadratic points. For \(k=3,\ q=2\), the ratio before the unique
quadratic is \(2<5=8-4+1\), giving a drop. The prefix after that
quadratic has weights \(1,1,1/3\) and \(R_2=7/5>1\).
The two cubic points give the subsequent rise. All cases are covered.
\end{proof}

This is a complete statement for all fields and all ranks, with no
unresolved large-rank range or omitted small field. It differs from
the rational-prime classification, whose unimodal ranks are
\(1,2,3\). Its proof uses exact finite-field factor counts and the
coefficient step identity, rather than a new zero-free estimate.

\section{How many modes: an exact large-field limit}

A positive peak plateau means a maximal run of equal positive entries
preceded by a strict increase and followed by a strict decrease.
Endpoint peaks are included. This convention ignores the initial
zero run and counts a two-entry flat maximum as one peak.

\begin{theorem}[Harmonic mode count]\label{thm:modeCount}
Fix \(k\geq3\), put \(r=k-1\), and let \(m_k\) be the unique integer with
\[
                       H_{m_k-1}<r<H_{m_k}.
\]
For all sufficiently large prime powers \(q\), the sequence
\(d_{q,k}\) has exactly \(m_k\) positive peak plateaus.
There is one at the last irreducible of each degree below \(m_k\),
and one in degree \(m_k\). If \(j\) is the position of this last peak
within its degree block, then
\[
 j=q^{m_k}(r-H_{m_k-1})+O_k(q^{m_k-1}).             \tag{17}
\]
Thus \(m_3=4,\ m_4=11\), and \(m_k\sim e^{k-1-\gamma}\) as \(k\to\infty\).
The order of quantifiers is fixed \(k\), then \(q\to\infty\);
no simultaneous growing-rank estimate is asserted.
\end{theorem}
\begin{proof}
For \(s\geq2\), \(H_s\) is not an integer: among \(1,\ldots,s\)
the denominator with largest power of two occurs exactly once,
so that power survives in the reduced denominator of the harmonic sum.
Harmonic divergence gives existence and uniqueness of \(m_k\).

The coefficient ratio \(R_r\) strictly decreases as a positive weight
is appended. Indeed its difference has the sign of
\(e_{r-2}e_r-e_{r-1}^2<0\), by Newton's inequalities; for \(r=1\)
use \(R_1=1/e_1\). Thus there is at most one peak plateau in a degree
block, and once \(R_r<1\), no later step can be an ascent.

Here are explicit bounds sufficient for the required approximation.
If \(\mu=\sum w_i\) and \(\nu=\sum w_i^2\), expansion of \(\mu^r\)
and a union bound over colliding positions give
\[
 0\leq\mu^r-r!e_r\leq\binom r2\nu\mu^{r-2}.
\]
For all irreducible prefixes,
\[
 \nu\leq\sum_{d\geq1}\frac{I_q(d)}{(q^d-1)^2}
       \leq4\sum_{d\geq1}\frac{q^{-d}}d\leq8/q.
\]
For fixed \(r\) and \(\mu\) in a fixed positive compact interval,
these inequalities imply \(R_r=r/\mu+O_r(1/q)\).
For each fixed degree \(n\),
\(I_q(n)=q^n/n+O_n(q^{n/2})\) when \(n\geq2\), while \(I_q(1)=q\).
Consequently the prefix after \(j\) points of degree \(n\) has
\[
 \mu=H_{n-1}+j/q^n+O_n(1/q)
\]
uniformly over that block.

At the end of degree \(m_k-1\) this gives \(R_r>1\);
monotonicity makes every earlier supported within-block step an
ascent. At the beginning and end of degree \(m_k\), the ratios
lie respectively above and below one. Monotonicity gives exactly
one positive peak plateau in that block and no later ascent.

For \(q\geq2r\), every degree boundary is a strict drop once the
density is positive: \(M_n\geq q\geq2r\), so (14) gives
\(R_r\leq q^n-1\) before the last point of degree \(n\).
This proves the stated number of peaks. At the final crossing,
\(R_r=1+O_k(q^{-m_k})\), since adjacent prefix means differ by
\(O(q^{-m_k})\); the same bound follows directly from the ratio
recursion and the uniform positive lower bound for \(e_r\).
Thus \(\mu=r+O_k(1/q)\), yielding (17); a one-position ambiguity
at a flat peak is absorbed by the error. Finally
\(H_n=\log n+\gamma+O(1/n)\) gives the asymptotic for \(m_k\).
\end{proof}

\paragraph{The constellation transfer must preserve the right information.}
A polynomial shift \(P\mapsto P+h\) with \(\deg h<\deg P\) preserves norm.
Any permutation made of such same-degree cycles leaves every univariate
norm-only prefix coefficient unchanged, because \(Q_{\sigma(i)}=Q_i\).
This observation concerns transported profiles; ordinary within-block
ascents still satisfy (13). The integer constellation theorem explicitly
has conclusion RH \emph{and} recurrence of the prescribed shape.
It does not assert that RH supplies prime-tuple recurrence.

\subsection{A character-sensitive family and its exact reference clock}

Fix odd \(q\) and a monic, nonconstant, squarefree
\(D\in\F_q[T]\) of degree \(d\). Extend the quadratic symbol by
\(\chi_D(P)=0\) at \(P\mid D\). Uniformly sample a monic polynomial
\(a\) of degree \(M\), and factor \(a^2-D\). Its local divisibility
probability is
\[
 \eta_D(P)=\frac{1+\chi_D(P)}{|P|}.
\]
Indeed, there are two, zero, or one roots of \(a^2=D\) modulo \(P\),
according as \(P\) splits, is inert, or divides \(D\). Thus ramified
factors have probability \(1/|P|\); split factors have probability
\(2/|P|\). These lie strictly below one because \(q\geq3\).

For every finite set of distinct \(P\)'s the divisibility indicators
are exactly independent once \(M\geq\sum_P\deg P\), by CRT. This is a
fixed-cutoff statement; no simultaneous growing-cutoff limit is used.
On any ordered list of positive-probability factors put
\(v_i=\eta_i/(1-\eta_i)\). Their \(k\)-th-factor density is exactly
\[
 b_{D,k}(i)=\eta_i\prod_{j<i}(1-\eta_j)
                        e_{k-1}(v_1,\ldots,v_{i-1}).
\]
When \(i\geq k\), with \(R=e_{k-2}/e_{k-1}\) on the preceding prefix,
\[
 \frac{b_{D,k}(i+1)}{b_{D,k}(i)}
  =\frac{\eta_{i+1}}{\eta_i}(1-\eta_i+\eta_i R).       \tag{18}
\]
If ramified factors are excluded and only split factors are indexed,
this specializes to
\[
 \frac{b_{D,k}(i+1)}{b_{D,k}(i)}
       =\frac{Q_i-2+2R}{Q_{i+1}},\qquad
 b_{D,k}(i+1)>b_{D,k}(i)
 \ \Longleftrightarrow\
 R>\frac{Q_{i+1}-Q_i+2}{2}.                         \tag{19}
\]
The ramified factors must either be retained with (18), or explicitly
excluded as in (19).

\begin{proposition}[Finite-field RH supplies the marked clock]
\label{prop:markedClock}
Let \(L_D(u)=\prod_P(1-\chi_D(P)u^{\deg P})^{-1}\).
This is a polynomial of degree \(A=d-1\); write
\(L_D(u)=\prod_{j=1}^{A}(1-\beta_j u)\).
Every \(|\beta_j|\leq\sqrt q\). Define the absolutely convergent constant
\[
 C_D=\log L_D(q^{-1})-
       \sum_P\sum_{\ell\geq2}
         \frac{1+\chi_D(P)^\ell}{\ell |P|^\ell}.
\]
For every integer \(n\geq1\),
\[
 \left|\sum_{\deg P\leq n}\eta_D(P)-H_n-C_D\right|
 \leq \frac{(A+4)q^{-(n+1)/2}}
              {(n+1)(1-q^{-1/2})}.                \tag{20}
\]
All Frobenius phases and all higher prime powers occur in the exact
signed identity
\[
\begin{split}
 \sum_{\deg P\leq n}\eta_D(P)-H_n
 ={}&-\sum_{m=1}^n\frac{\sum_{j=1}^{A}\beta_j^m}{m q^m}\\
 &-\sum_{m=1}^n\frac1{m q^m}
     \sum_{\substack{e\mid m\\e<m}}e
          \sum_{\deg P=e}(1+\chi_D(P)^{m/e}).
\end{split}\tag{21}
\]
\end{proposition}
\begin{proof}
The smooth projective model of \(y^2=D(T)\) has genus
\(\lfloor(d-1)/2\rfloor\). If \(d\) is odd it has one rational
point at infinity; if \(d\) is even, monicity gives two. Its point-count
formula and the Euler product show that \(L_D\) is the curve's
zeta numerator when \(d\) is odd, and \((1-u)\) times that numerator
when \(d\) is even. Thus its reciprocal roots are the curve's
Frobenius roots, and, in the even case, one additional root \(1\).
This proves the asserted degree and root bound using
finite-field RH~\cite{Milne}. The statement includes \(d=1\), for which
\(L_D=1\) and the reciprocal-root sum is empty.

For \(S_m=\sum_{\deg P=m}\chi_D(P)\), logarithmic differentiation gives
\[
 \sum_{e\mid m}e\sum_{\deg P=e}\chi_D(P)^{m/e}
      =-\sum_j\beta_j^m.
\]
Add \(q^m=\sum_{e\mid m}eI_q(e)\), isolate the \(e=m\) terms,
and divide by \(m q^m\). Since
\(\sum_{\deg P=m}\eta_D(P)=q^{-m}(I_q(m)+S_m)\),
summing proves (21), with no omitted compensation or prime-power terms.

The proper-divisor contribution to each of the two degree identities
has magnitude at most
\(\sum_{e\leq m/2}q^e\leq2q^{m/2}\).
Consequently
\[
 \left|q^{-m}(I_q(m)+S_m)-\frac1m\right|
       \leq\frac{(A+4)q^{-m/2}}m.
\]
For \(m=1\) the proper-divisor sums are empty and the same bound holds.
This proves absolute convergence of the degree errors. In (21), the
first infinite sum is \(\log L_D(q^{-1})\), defined by its convergent
real power series. The second infinite sum, after writing \(m=e\ell\),
is the displayed prime-power sum defining \(C_D\).
The latter is absolutely convergent since its summands are bounded
by \(2/(\ell |P|^\ell)\), \(\ell\geq2\).
Finally the tail of the degree errors is bounded by
\[
 (A+4)\sum_{m>n}\frac{q^{-m/2}}m
 \leq\frac{(A+4)q^{-(n+1)/2}}{(n+1)(1-q^{-1/2})},
\]
which proves (20).
\end{proof}

For the split-only profile, subtract
\(\sum_{P\mid D,\ \deg P\leq n}|P|^{-1}\) from the clock.
For \(n\geq d\) this simply replaces \(C_D\) by
\(C_D-\sum_{P\mid D}|P|^{-1}\).
To obtain the coefficient clock of (18), use the exact odds identity
\[
 \sum_{\deg P\leq n}\frac{\eta_D(P)}{1-\eta_D(P)}
  =\sum_{\deg P\leq n}\eta_D(P)
       +\sum_{\deg P\leq n}\frac{\eta_D(P)^2}{1-\eta_D(P)}.
\]
The second sum converges absolutely, with tail \(O_q(q^{-n}/n)\):
use \(\eta_D(P)\leq2q^{-\deg P}\), \(1-\eta_D(P)\geq1/3\),
and \(I_q(m)\leq q^m/m\).
Thus the actual character-sensitive coefficient mean has a
square-root-accurate harmonic reference as well.

This gives a concrete analytic input for transferring the ascent
comparison: the reference clock, every local sampling weight, the exact
step ratio, and a uniform error are available. Transport by a
same-degree prime cycle can now change the marked coefficients, since
it can exchange primes with different \(\chi_D\)-values.
Supply of the required shifted irreducible tuples remains a separate
hypothesis; the integer theorem also states RH and shape recurrence as
a conjunction. No tuple-count estimate is deduced from RH alone.

\section{Quadratic Galois profiles: finite words and sharp mode complexity}

Throughout this section \(q\) is a fixed odd prime power and
\(D\in\F_q[T]\) is monic, nonconstant and squarefree. Order the
unramified split primes of \(\F_q(T)(\sqrt D)\) by degree, breaking
ties arbitrarily. Write \(b_{D,k}\) for the limiting density of the
\(k\)-th distinct prime on this list dividing \(a^2-D\), where \(a\)
is sampled uniformly among monic polynomials of degree \(M\) and
\(M\to\infty\) at each fixed prime cutoff. Ramified factors are
excluded from this rank. The local probabilities are \(2/|P|\),
so (19) applies. These are actual arithmetic sampling profiles,
not freely assigned Bernoulli weights.

\begin{lemma}[Squarefree realization of prescribed splitting]
\label{lem:splittingPrescription}
Let \(U\) be a nonconstant monic squarefree polynomial of degree \(m\),
and let \(c\) be a unit modulo \(U\). There exists a monic squarefree
\(D\equiv c\pmod U\) of degree \(3m+2\).
Consequently any assignment of split or inert status to finitely many
primes can be realized by a geometric quadratic extension.
\end{lemma}
\begin{proof}
There are \(q^{\ell-m}\) monic degree-\(\ell\) representatives of the
class, where \(\ell=3m+2\). No prime dividing \(U\) divides a representative.
For a prime \(P\nmid U\) of degree \(e\leq m+1\), CRT says the fraction
divisible by \(P^2\) is \(q^{-2e}\). The union bound for these primes is
less than
\(\sum_{e\geq1}q^{-e}=1/(q-1)\), since \(I_q(e)\leq q^e\).
For \(m+2\leq e\leq\ell/2\), ignore the congruence condition:
there are at most \(I_q(e)q^{\ell-2e}\leq q^{\ell-e}\)
monic degree-\(\ell\) polynomials divisible by such a square.
Their fraction in our class is at most
\[
 q^m\sum_{e\geq m+2}q^{-e}=\frac1{q(q-1)}.
\]
The total is less than
\((q+1)/(q(q-1))\leq2/3<1\). A squarefree representative therefore exists.

For the last assertion take \(U\) to be the product of the prescribed
primes. At each prime prescribe a nonzero square residue for splitting,
or a nonsquare residue for inertia; both exist in its odd residue field.
CRT produces a unit \(c\), and the first part gives \(D\).
Its positive-degree squarefreeness makes \(D\) nonsquare even over
\(\overline{\F}_q(T)\), so the geometric and arithmetic Galois groups
are both the two-element group.
\end{proof}

\begin{theorem}[Every finite ascent word occurs]\label{thm:quadraticWords}
Fix \(q\) odd, \(k\geq2\), and a finite word
\(\epsilon_1,\ldots,\epsilon_t\in\{+,-\}\).
There is a monic squarefree \(D\) for which a contiguous segment of the
positive profile \(b_{D,k}\) has exactly these strict comparison signs.
The field \(q\), rank \(k\), and quadratic Galois group stay fixed.
\end{theorem}
\begin{proof}
Put \(r=k-1\) and \(s=2r+1+t\). Choose \(N\) large enough that every
degree \(n\in[N,N+t]\) has at least \(s\) primes, and
\(2s/(q^N-2)<r\). Such \(N\) exists by the irreducible count formula;
\(N\geq\max(4,4s)\) suffices using
\(I_q(n)\geq q^n/(2n)\geq n/2\) for \(n\geq4\).
The first inequality follows from the proper-subfield bound above,
and the second from \(q^n\geq3^n\geq n^2\).

Start with \(2r+1\) selected split primes of degree \(N\).
Read the word in order. For a plus, choose the next selected prime
in the same degree; for a minus, choose it in the next degree.
Declare all other primes through degree \(N+t\) inert.
Lemma~\ref{lem:splittingPrescription} realizes this finite prescription.
There are \(s\) selected primes, and the sum of all their odds
\(w=2/(q^{\deg P}-2)\) is less than \(r\).

For every supported comparison in the prescribed word,
(14) gives \(R_r\geq r/\mu>1\); equal-degree steps are therefore
strict ascents by (19). At a degree jump the preceding prefix has
at least \(2r\) weights. If the current norm is \(Q\), (14) gives
\[
 R_r\leq\frac{r(Q-2)}{2(m-r+1)}
       \leq\frac{r(Q-2)}{2(r+1)}
       <\frac{Q-2}{2},
\]
where \(m\geq2r\) is the prefix size. This is below the threshold
\(((q-1)Q+2)/2\), so the jump is a strict descent.
The initial \(2r+1\) entries ensure positive density throughout the
recorded segment. Every letter is realized, completing the proof.
\end{proof}

Let \(\mathcal N_{q,k}(D)\) count the positive peak plateaus of
\(b_{D,k}\), with the convention used above. For integers \(H\geq1\),
put
\[
 \mathcal N_{q,k}^{*}(H)=
 \max_{\substack{D\ {\rm monic,\ squarefree}\\1\leq\deg D\leq H}}
                    \mathcal N_{q,k}(D).
\]
The maximum is over a finite, nonempty set.

\begin{theorem}[Sharp logarithmic order in the conductor degree]
\label{thm:quadraticModeComplexity}
For fixed odd \(q\), one has \(\mathcal N_{q,1}^{*}(H)=1\).
For every fixed \(k\geq2\),
\[
                    \mathcal N_{q,k}^{*}(H)
                         \asymp_{q,k}\log H
                    \qquad(H\to\infty).          \tag{22}
\]
In particular, each quadratic profile has finitely many peaks, but
there is no bound depending only on \(q,k\) and its two-element
Galois group.

More explicitly, let \(r=k-1\) and
\(C_{q,k}=3q^{4r+7}/(q-1)+2\). For
\(H\geq C_{q,k}q^7\),
\[
 \mathcal N_{q,k}^{*}(H)\geq
 \left\lfloor\frac{\log_q H-\log_q C_{q,k}}7\right\rfloor.   \tag{23}
\]
For every individual \(D\) of degree \(d\geq1\), set
\[
 L=\left\lceil2\log_q(4(d+3))\right\rceil,\qquad
 B=\left\lceil e^{4r+2}L\right\rceil.
\]
Then
\[
                         \mathcal N_{q,k}(D)\leq B.        \tag{24}
\]
\end{theorem}
\begin{proof}
For \(k=1\), the split-prime step ratio is
\((Q_i-2)/Q_{i+1}<1\), so the first entry is the unique endpoint peak.

For the upper bound write \(s_D(n)\) for the number of unramified
split primes of degree \(n\), \(j_D(n)\) for the number dividing \(D\),
and \(S_n=\sum_{\deg P=n}\chi_D(P)\). Then
\[
 2s_D(n)=I_q(n)+S_n-j_D(n),\qquad j_D(n)\leq d/n.
\]
The logarithmic-derivative argument in
Proposition~\ref{prop:markedClock} gives
\[
 \left|I_q(n)+S_n-\frac{q^n}{n}\right|
       \leq\frac{(d+3)q^{n/2}}n .
\]
For every \(n\geq L\), therefore,
\[
 2s_D(n)\geq\frac{q^n}{n}
       \bigl(1-(d+3)q^{-n/2}-d q^{-n}\bigr)
       \geq\frac{q^n}{2n}.
\]
Here \(q^{n/2}\geq4(d+3)\), so the two subtractions are at most
\(1/4\) and \(1/16\). In particular there are split primes in each
such degree. The total odds of split primes in degrees \(L,\ldots,B\)
are at least
\[
 \sum_{n=L}^B\frac{2s_D(n)}{q^n}
 \geq\frac12\sum_{n=L}^B\frac1n
 \geq\frac12\log\frac{B+1}{L}\geq2r+1.
\]
Each of these odds is at most one, since
\(q^L\geq16(d+3)^2\geq256\).
In the prefix before the last degree-\(B\) split prime, their total
is still at least \(2r\).

For any positive weights of mean sum \(\mu\) and largest weight
\(w_{\max}\), the coefficient identity used in (14) also gives
\[
 \frac{e_{r-1}}{e_r}
       \leq\frac{r}{\mu-(r-1)w_{\max}}
 \quad\hbox{when the denominator is positive}.             \tag{25}
\]
Indeed, after choosing \(r-1\) weights their removal leaves sum
at least \(\mu-(r-1)w_{\max}\).
Applied just to the tail weights in that prefix, (25) is at most
\(r/(r+1)<1\). Appending the earlier weights can only decrease
the ratio, by Newton's inequality as in the mode-count proof.
Thus the outgoing step from the last degree-\(B\) prime, and every
later step, is strictly decreasing by (19).

Within each degree block, consecutive step ratios strictly decrease
because \(R_r\) strictly decreases. A block contributes at most one
positive peak plateau; the initial zero run contributes none.
All peaks have degree at most \(B\), proving (24).
Since \(B=O_{q,k}(\log(d+2))\), this proves the upper order in (22).

For the lower bound fix any desired number \(v\geq1\) of strict peaks.
Put
\[
 N=\max(4,4r,6(v+1)),\qquad b=N+v.
\]
Prescribe every prime of degree less than \(N\) inert and every
prime of degree \(N,\ldots,b\) split.
For \(n\geq N\), \(I_q(n)\geq n/2\geq2r\).
The total odds in the prescribed split range are bounded by
\[
 \sum_{n=N}^{b}\frac{2I_q(n)}{q^n-2}
       \leq\sum_{n=N}^{b}\frac3n
       \leq\frac{3(v+1)}N\leq\frac12<r.
\]
Every supported within-block step is therefore a strict ascent.
Before the last prime of each degree \(n=N,\ldots,b-1\), the
prefix has at least \(I_q(n)-1\geq2r-1\) weights, and (14) gives
\(R_r\leq(q^n-2)/2\). The next degree is a strict descent by (19).
Hence the last prime of each of these \(v\) blocks is a strict peak.

Apply Lemma~\ref{lem:splittingPrescription} with the product \(U\)
of all primes through degree \(b\). Since
\[
 \deg U=\sum_{n\leq b}nI_q(n)
     \leq\sum_{n\leq b}q^n<\frac{q^{b+1}}{q-1},
 \qquad b\leq4r+7v+6,
\]
the realizing squarefree \(D\) has
\[
 \deg D=3\deg U+2
       \leq C_{q,k}q^{7v}.
\]
Taking
\(v=\lfloor(\log_q H-\log_q C_{q,k})/7\rfloor\)
proves (23) and the lower order in (22).
\end{proof}

\paragraph{What the Galois connection contributes.}
The group determines the possible local types, but the conductor can
program a long initial pattern of those types. The coefficient comparisons
turn that arithmetic pattern into actual ascents and descents.
Finite-field RH then controls how long this freedom can persist:
the maximal peak count is logarithmic in \(\deg D\), the degree
of the finite quadratic conductor. The lower bound is an elementary
CRT and squarefree-counting construction; the upper bound uses
the proved finite-field RH character estimate. Both directions concern
fixed \(q\), not a large-field limit.

\paragraph{Scope, closest literature, and formalization status.}
The existing unmarked classification and fixed-rank large-\(q\) mode
law above do not imply these fixed-\(q\), varying-conductor statements.
Cambie~\cite{Cambie} studies integer divisor and prime-factor profiles,
including many extrema for the interval-divisor statistic.
Jung~\cite{Jung} supplies function-field Mertens asymptotics in residue
classes. Neither statement is the finite-word realization or the sharp
logarithmic peak complexity of the quadratic sampling profile considered
here. A targeted search of the main repository's public and private
plans found no matching unfinished target; a complete cross-repository
priority audit has not been performed. No matching theorem was identified
in the literature search, which does not establish priority.

These are ordinary proofs, not newly compiled Lean declarations.
The concrete formalization gaps are the finite-field CRT sampling adapter,
the quadratic character degree-count estimate from the curve zeta formula,
and the squarefree residue-class realization lemma. The numerical companion
\texttt{quadratic\_profile\_word\_probe.py} independently checks the
coefficient consumer for \(q\in\{3,5,7\}\), \(2\leq k\leq6\),
and every sign word of length at most six: 1905 words and 9630 signed
steps pass in exact rational arithmetic. It does not construct the
large CRT polynomial or replace any of the universal proofs.

\section{Transferring the H1-dependent ascent scale}
\label{sec:gapScaleTransfer}

The relevant integer input is the maintained exact critical-threshold
theorem, not merely an analogy with twin primes. If
\(H_1=\liminf(p_{j+1}-p_j)\), its ordinary canonical late-ascent
coefficient is
\[
                         \beta_{\mathbf Z}=\frac1{H_1+1}.
                                                               \tag{29}
\]
The maintained moving-rank RH criterion compares the actual coefficient
ratio with its canonical reference at every fixed positive rank-to-clock
ratio. These are two connected statements, and their arithmetic
hypotheses must be transported separately. The sources are
\texttt{GapRecurrence.lean}, \texttt{PrimeProfileMinimax.lean}, and
the moving-rank criterion in the maintained integer paper.
Equation (29) states the threshold property: smaller coefficients have
unbounded late-ascent supply, while larger coefficients force terminal
descent. No endpoint assertion or new integer gap bound is assumed.

\subsection{The actual probability law on a curve}
Let \(C/\F_q\) be as above and enumerate its closed points
\(P_1,P_2,\ldots\) in nondecreasing degree. Put \(Q_i=q^{\deg P_i}\).
Sample uniformly from the effective divisors of degree \(M\), and
then let \(M\to\infty\) at a fixed finite closed-point cutoff.
Let \(d_{C,k}(i)\) denote the limiting probability that \(P_i\) is
the \(k\)-th distinct point in the divisor's support.

This is a genuine sampling law. If \(a_M\) counts effective divisors
of degree \(M\), the curve zeta identity gives
\[
 \sum_{M\geq0}a_M u^M
       =\frac{P_C(u)}{(1-u)(1-qu)}.
\]
For \(M\geq2g\), coefficient extraction yields the exact expression
\[
 a_M=\frac{q^{M+1}P_C(1/q)-P_C(1)}{q-1}.
\]
The leading constant is positive: in the conjugate-paired product
\(P_C(1/q)=\prod_i(1-\alpha_i/q)\), each conjugate pair is positive,
and any real reciprocal root is \(\pm\sqrt q\).
For any fixed effective divisor \(E\),
\(\Pr(E\leq\mathcal D_M)=a_{M-\deg E}/a_M\to q^{-\deg E}\).
Inclusion-exclusion therefore gives independent limiting divisibility
indicators on every fixed finite set of points.
There is no simultaneous growing-cutoff limit here.

Consequently, with \(w_i=(Q_i-1)^{-1}\),
\[
 d_{C,k}(i)=\frac1{Q_i}\prod_{j<i}(1-Q_j^{-1})e_{k-1}(w_j:j<i).
\]
For \(i\geq k\geq2\), writing \(R_r=e_{r-1}/e_r,\ r=k-1\),
\[
 \frac{d_{C,k}(i+1)}{d_{C,k}(i)}
      =\frac{Q_i-1+R_r(i)}{Q_{i+1}}.               \tag{30}
\]
In \(\F_q[T]\), sampling monic polynomials gives the affine version
already used above. The finitely many points at infinity change
constants but not the critical coefficient below.

\begin{theorem}[Fixed-field, growing-rank critical scale]
\label{thm:fixedFieldCritical}
Fix \(C/\F_q\) and any ordering within its degree blocks.
For every \(0<b<1\) and all sufficiently large \(k\), there is
an ascent \(d_{C,k}(i)<d_{C,k}(i+1)\) with
\[
                         b k\leq\log\log Q_i.
\]
Their degrees tend to infinity as \(k\to\infty\).
For every \(b>1\), all sufficiently large \(k\) satisfy
\[
 b k\leq\log\log Q_i
       \quad\Longrightarrow\quad
 d_{C,k}(i+1)<d_{C,k}(i).
\]
The rank and point cutoffs in this assertion are uniform.
Thus the critical coefficient is exactly \(1\).
For all sufficiently large \(k\), a last ascent exists; if \(Q_k^{\rm last}\)
is its starting norm, then
\[
                 \lim_{k\to\infty}
                    \frac{\log\log Q_k^{\rm last}}k=1.          \tag{31}
\]
This is a fixed-field, growing-rank theorem, separate from the earlier
fixed-rank, growing-field mode count.
\end{theorem}
\begin{proof}
Weil's point count and the exact closed-point inversion give
\(B_n=q^n/n+O_{C,q}(q^{n/2}/n)\).
It follows that
\[
 \mu_n:=\sum_{d\leq n}\frac{B_d}{q^d-1}
       =H_n+\kappa_C+O_{C,q}(q^{-n/2}/n)
       =\log n+c_C+O_{C,q}(1/n).                  \tag{32}
\]
Indeed each summand minus \(1/d\) is \(O_{C,q}(q^{-d/2}/d)\),
so its error series converges absolutely.
Every prefix before a degree-\(n\) point lies between the full
degree-\((n-1)\) and degree-\(n\) prefixes. Their mean difference
is \(O_{C,q}(1/n)\), so (32) is uniform within each block.
Also \(B_n\geq2\) for all sufficiently large \(n\).

Fix \(0<b<a<1\). For each large \(k\), take
\(n=\lceil e^{ak}/\log q\rceil\).
Then \(\log\log(q^n)=ak+o(1)\geq bk\), while
\(\mu_{n-1}=ak+O_{C,q}(1)<k-1\).
There are more than \(k\) points in earlier degrees.
Before the first point of degree \(n\), (14) gives
\(R_{k-1}\geq(k-1)/\mu_{n-1}>1\).
The first two points in that block therefore form an ascent by (30).
This proves the lower statement for every sufficiently large rank.

For the upper statement fix \(b>1\), put
\(\delta=(b-1)/4\), and remove all points through a fixed degree \(D_0\)
large enough that every remaining weight is at most \(\delta\).
If \(bk\leq\log\log(q^n)\), (32) gives a tail-prefix mean
at least \(bk-A\), for a constant \(A=A(C,q,b)\) independent
of \(k,n\) and of position inside the block.
For all sufficiently large \(k\), this exceeds
\((k-1)+(k-2)\delta\), since \(b-\delta>1\).
There are at least \(k-1\) tail weights; otherwise their total
would be at most \((k-2)\delta\).
Apply (25) just to those weights:
\[
 R_{k-1}^{\rm tail}
       \leq\frac{k-1}{\mu_{\rm tail}-(k-2)\delta}<1.
\]
Appending the removed earlier weights can only decrease this ratio,
by Newton's inequality. Hence the actual \(R_{k-1}<1\).
Since \(Q_{i+1}\geq Q_i\), (30) forces strict descent.
This argument is uniform over all indices in the stated region
and does not use a fixed-rank coefficient approximation.

Taking any \(b>1\) leaves only finitely many possible ascent indices.
The lower statement supplies at least one for every large \(k\).
Applying the two statements for coefficients arbitrarily close to one
proves (31).
\end{proof}

\subsection{Which gap has the unmarked transfer recovered?}
The exact integer step compares \(R_r\) with \(p_{i+1}-p_i+1\).
Equation (30) compares it with
\[
                         Q_{i+1}-Q_i+1.
\]
Distinct closed points of the same degree have \(Q_{i+1}-Q_i=0\).
Consequently the least recurring \emph{norm difference} is zero,
and the transferred threshold is \(1/(0+1)=1\), exactly as the
theorem proves. This is not the polynomial distance
\(|P_{i+1}-P_i|=q^{\deg(P_{i+1}-P_i)}\).
A permutation within a degree block leaves the complete unmarked
probability profile unchanged, even when it changes the polynomial
differences of adjacent labels.

The distinction holds at every large degree, not merely in a finite
test. Over \(\F_2[T]\), no pair \(f,f+1\) of degree greater than one
can both be irreducible: one has constant coefficient zero, and is
divisible by \(T\). Nonetheless Theorem~\ref{thm:fixedFieldCritical}
gives critical coefficient \(1\) for the ordinary polynomial profile.
Over \(\F_7[T]\), by contrast,
\[
                         T^{3^m}-2,\qquad T^{3^m}-4
\]
are both irreducible for every \(m\geq0\), by the binomial
irreducibility criterion: \(2\) and \(4\) are noncubes in \(\F_7\);
see also Burrin--Issac~\cite{BurrinIssac}.
The ordinary critical coefficient is still \(1\).
The unmarked threshold therefore does not distinguish these
different geometric recurrence situations.

\begin{proposition}[The shift-marked recurrence transfer]
\label{prop:shiftMarkedTransfer}
Fix \(0\ne h\in\F_q[T]\) and \(0<b<1\).
The following are equivalent:
\begin{enumerate}
\item For every \(N\), there is a monic \(P\) with
\(\deg P\geq N\) such that \(P,P+h\) are irreducible.
\item There is a degree-compatible ordering \(P_1,P_2,\ldots\) of
the monic irreducibles such that for every \(K,N\), there are
\(k\geq K\) and \(i\) with
\[
 \deg P_i\geq N,\quad P_{i+1}=P_i+h,\quad
 bk\leq\log\log|P_i|,\quad
 d_{q,k}(i)<d_{q,k}(i+1).
\]
\end{enumerate}
For \(b>1\), the second property is false for every ordering.
\end{proposition}
\begin{proof}
The second property immediately supplies the pairs in the first.
Conversely, choose one pair in each of an unbounded sequence of distinct
degrees \(n>\deg h\), and place its two elements consecutively,
in the displayed orientation, within that degree block.
This defines one fixed global ordering.

Choose \(b<a<1\). For any such large degree \(n\), set
\(k=\lfloor\log\log(q^n)/a\rfloor\).
Then \(k\to\infty\), \(bk\leq\log\log(q^n)\), and the mean of
every prefix in that block is
\(\log n+O_q(1)<k-1\).
There are more than \(k\) earlier-degree primes.
Thus (14) gives \(R_{k-1}>1\) before the selected pair,
and (30) gives its strict ascent.
Arbitrarily large selected degrees accommodate both \(K\) and \(N\).
The upper statement follows from
Theorem~\ref{thm:fixedFieldCritical}, whose proof also applies to the
affine polynomial prime set.
\end{proof}

This proposition retains the geometric shift, rather than relabelling
zero norm difference as a polynomial gap. Its marked critical coefficient
is one whenever the selected shift recurs. We now transfer the canonical
moving-rank reference as well, retaining the zero signal and a genuine
shift response in the same positive probability model.

\subsection{Positive laws and the uniform reference comparison}
For local probabilities \(0<\eta_P<1\), put \(w_P=\eta_P/(1-\eta_P)\).
For a finite prime prefix \(S\), write \(e_j(S)=e_j((w_P)_{P\in S})\),
and \(R_r(S)=e_{r-1}(S)/e_r(S)\). This is an actual adjacent-rank
density ratio: at the next prime, the common factor
\(\eta_P\prod_{Q\in S}(1-\eta_Q)\) cancels from \(d_r/d_{r+1}\).
There is no division by a density at a prime with zero local probability.

Suppose \(\eta_P\ll|P|^{-1}\). Define
\[
 {\cal E}(z)=\prod_P(1+w_Pz)e^{-\eta_Pz},\qquad
 F_r(L)=[z^r]e^{Lz}{\cal E}(z),\qquad
 {\cal R}_r(L)=\frac{F_{r-1}(L)}{F_r(L)}.
\]
The product converges absolutely, locally uniformly, since
\(w_P-\eta_P=O(|P|^{-2})\) and \(\sum_P|P|^{-2}<\infty\).
It is strictly positive for real \(z\geq0\).
Absolute convergence makes it invariant under degree-compatible
permutations, so the original and swapped orders use the same reference.
Put \(B(S)=\sum_{P\in S}\eta_P\).
The maintained uniform coefficient comparison gives, on each fixed band
\(0<\alpha_0\leq r/L\leq\alpha_1<\infty\), when \(|B(S)-L|\leq1\),
\begin{align}
 R_r(S)-{\cal R}_r(L)
 &=-\frac r{L^2}(B(S)-L)
   +O\left(\frac{|B(S)-L|}{L^2}+\frac1{r q^n}\right),
 \label{eq:markedLinearization}\\
 {\cal R}'_r(L)&=-\frac r{L^2}+O(L^{-2}).
 \label{eq:markedSlope}
\end{align}
Here \(S\) consists of the primes through degree \(n\), with at most two
specified primes removed. Constants may depend on the law, \(q\), and the
band, but not on \(n,r,L\). All ranks in this subsection tend to infinity.

For completeness, the application of the integer-indexed provider is
legitimate without assuming integer RH. Enumerate the polynomial primes
by degree and label the \(j\)-th one by the \(j\)-th rational prime.
The prime polynomial formula and the ordinary prime number theorem give
\(p_j\asymp_q |P_j|\). Thus \(w_{P_j}\ll1/(p_j-1)\),
\(\eta_{P_j}\ll1/(p_j-1)\), and
\(|w_{P_j}-\eta_{P_j}|\ll1/(p_j-1)^2\), exactly the hypotheses of
\texttt{PrimeClockLaw}. A prefix ending in degree \(n\) has rational cutoff
comparable to \(q^n\). Removing finitely many terminal weights is also
handled by their exact coefficient identities. The reference positivity,
slope bound, and error in \eqref{eq:markedLinearization} are precisely the
maintained providers \texttt{exists\_densityRatio\_linearization} and
\texttt{exists\_profile\_reference\_slope\_bound}.

The probability limit here always takes the sampling degree to infinity
at a fixed finite divisor-prime cutoff first. The growing variable \(n\)
is the divisor cutoff degree, not a substitution into the sampled family.

\begin{lemma}[Coupled replacement response]\label{lem:coupledResponse}
On a common background put \(e_0=e_r\), \(e_1=e_{r-1}\), and
\(e_2=e_{r-2}\). Inserting a weight \(a\) or \(b\) gives
\begin{equation}\label{eq:exactCoupledResponse}
 \frac{e_1+be_2}{e_0+be_1}
 -\frac{e_1+ae_2}{e_0+ae_1}
 =\frac{(a-b)(e_1^2-e_0e_2)}
 {(e_0+ae_1)(e_0+be_1)}.
\end{equation}
If the background has at least \(r\) positive weights, \(r\geq2\),
and \(a>b\geq0\), this response is strictly positive.
If \(r/L\) stays in a fixed positive compact band, the clock is
\(L+O(1)\), and \(a,b=O(q^{-n})\), then
\begin{equation}\label{eq:coupledAsymptotic}
 R_r(S\cup\{b\})-R_r(S\cup\{a\})
 =(a-b)\frac r{L^2}(1+o(1)).
\end{equation}
\end{lemma}
\begin{proof}
Cross multiplication proves \eqref{eq:exactCoupledResponse}, including
every background cross term. Newton's strict inequality gives
\(e_1^2-e_0e_2>0\).

Here are the error estimates needed for the asymptotic, rather than a
subtraction of the two separate reference errors. Let \(U\) be the
arithmetic clock of the full prefix. The exact identities
\(F'_r=F_{r-1}\) give
\[
 {\cal R}'_r(U)=\frac{F_{r-2}(U)}{F_r(U)}
                     -{\cal R}_r(U)^2.
\]
At its own clock \(U\), the finite-prefix ratios of ranks \(r,r-1\)
differ from their infinite-reference ratios by \(O(1/(r q^n))\).
These ratios are uniformly bounded on the entire claimed band.
To check the denominator in (25), remove a fixed finite prefix until
each remaining odds weight is at most
\(\delta=1/(4(1+\alpha_1))\).
The tail odds sum is \(L+O(1)\), because
\(\sum_P(w_P-\eta_P)\) converges. Thus
\(\mu_{\rm tail}-(r-1)\delta\geq L/2\) for all sufficiently large
\(L\). The tail has at least \(r\) weights, since otherwise its mean
would be at most \((r-1)\delta<L/2\).
Formula (25) bounds its ratios through rank \(r\) by \(2\alpha_1+1\).
Newton monotonicity bounds the full-prefix ratios by these tail
ratios. The same argument applies to the common background after
the two terminal weights are removed.
Apply \eqref{eq:markedSlope} to conclude that for the full prefix
\[
 \left(\frac{e_{r-1}}{e_r}\right)^2-\frac{e_{r-2}}{e_r}
       =\frac r{L^2}+O(L^{-2})+O(q^{-n}).
\]
Removing either of the last two weights changes each of the needed
bounded ratios by \(O(q^{-n})\): this follows directly by writing
\(e_j(S\cup\{a\})=e_j(S)+a e_{j-1}(S)\), with all denominators
positive. Use the uniform background bounds just proved; its clock
differs by \(O(q^{-n})\). Therefore the same determinant estimate
holds on the common background. Divide \eqref{eq:exactCoupledResponse}
by \(e_0^2\). Its two denominator factors become
\(1+a e_1/e_0\) and \(1+b e_1/e_0\), both \(1+O(q^{-n})\).
The relative error is
\(O(1/r)+O(L^2/(r q^n))=o(1)\). This proves the assertion with the
common zero-clock error still coupled.

The uniformity over the terminal pair does not assume a new
family-uniform theorem for arbitrarily relabelled laws. Apply the
maintained provider at full degree cutoffs in one fixed ordering.
For a single deleted weight \(w=\eta/(1-\eta)=O(q^{-n})\), the exact
response with the replacement weight zero gives
\[
 R_r(S_n\setminus\{P\})-R_r(S_n)
       =\eta\,\frac r{L^2}
          +O\left(\frac1{q^nL^2}+q^{-2n}\right).
\]
The slope estimate and the mean value theorem give, for the same rank,
\[
 {\cal R}_r(L)-{\cal R}_r(L+\eta)
       =\eta\,\frac r{L^2}+O(1/(q^nL^2)).
\]
Apply the full-prefix comparison at reference clock \(L+\eta\).
The two displayed leading terms cancel, and
\(B(S_n\setminus\{P\})-L=B(S_n)-(L+\eta)\) exactly.
The remaining deletion error is \(O(1/(r q^n))\), since
\(r\asymp L=O(\log n)\). Repeating at most twice establishes
\eqref{eq:markedLinearization} for all the partial prefixes used here
with constants independent of which terminal primes were removed.
\end{proof}

\begin{lemma}[Signed packet stability from only two coordinates]
\label{lem:packetSignedStability}
Let \(a\ne b\) belong to a finite class set \(G\), and write
\(v_c=\mathbf1_{c=a}-\mathbf1_{c=b}\).
For arbitrary real errors \(g_c\), retaining every cross term gives
\[
 \sum_{c\in G}(v_c+g_c)^2
    =2+2(g_a-g_b)+\sum_{c\in G}g_c^2.
\]
In particular, the two one-sided bounds \(g_a\geq-\varepsilon\) and
\(g_b\leq\varepsilon\) imply
\[
                  \sum_{c\in G}(v_c+g_c)^2\geq2-4\varepsilon.
\]
There is no hypothesis on any other error coordinate and no loss
proportional to the packet size in this lower bound.
\end{lemma}
\begin{proof}
Expand the square before estimating. The exact identities
\(\sum_c v_c^2=2\) and \(\sum_c g_cv_c=g_a-g_b\) give the formula.
Now \(g_a-g_b\geq-2\varepsilon\), and every retained error square
is nonnegative. The three identities and this lower bound are
compiled as \texttt{class\_packet\_unit\_energy},
\texttt{class\_packet\_cross\_sum},
\texttt{class\_packet\_energy\_expansion}, and
\texttt{class\_packet\_energy\_lower}.
\end{proof}

\subsection{A complete finite packet with no shift blind spot}
Fix \(0\ne h\in\F_q[T]\). Choose a monic irreducible modulus
\(\mathfrak r\nmid h\) of degree \(d\geq2\), and put
\[
 G=(\F_q[T]/\mathfrak r)^\times,\qquad
 \phi=q^d-1,\qquad \kappa=1+\phi^{-1}.
\]
For a polynomial \(b\), let
\(b_c=\prod_{P\bmod\mathfrak r=c}P^{v_P(b)}\), for each \(c\in G\).
Sample independent uniform monic \(a,b\) of degree \(M\), and use the
computable arithmetic family \(F_c(a,b)=a b_c\). This family uses a
factorization filter; it is not asserted to be a fixed polynomial
expression in \(a,b\). At each fixed finite cutoff CRT gives independent
local divisibility indicators, exactly once \(M\) exceeds the degree
of the product of the cutoff primes, with
\begin{equation}\label{eq:classCarrierProbability}
 \eta_c(P)=\frac1{|P|}
       +\mathbf1_{[P]=c}\left(\frac1{|P|}-\frac1{|P|^2}\right).
\end{equation}
The indicator is zero at \(P=\mathfrak r\).
Indeed, the extra event is \(P\mid b\) in the marked class, overlapping
\(P\mid a\) with probability \(|P|^{-2}\). In particular
\(0<\eta_c(P)<1\) at every prime, including the small and ramified ones.

\begin{theorem}[Zero-preserving marked reference packet]
\label{thm:completeMarkedPacket}
For these \(\phi\) positive laws the following assertions hold.
\begin{enumerate}
\item Their clocks have constants \(C_c\) such that
\[
 B_c(n)=\sum_{\deg P\leq n}\eta_c(P)
       =\kappa H_n+C_c+E_c(n),\qquad
 E_c(n)=O_{q,d}(q^{-n/2}/n),
\]
uniformly over the finite packet.
After explicitly removing the prime-power, ramification, and overlap
tails from \(E_c(n)\), its finite Fourier transform recovers every
nonprincipal Dirichlet \(L(u,\chi)\) modulo \(\mathfrak r\), including
all of its inverse roots and their multiplicities.
\item Suppose \(P,Q=P+h\) are actual irreducibles of degree \(n>d\).
Place them last in that degree block, \(P\) followed by \(Q\), and
compare this prefix cut with the cut in the swapped ordering.
Use in each class
\[
 L_c(n)=\kappa H_n+C_c-\eta_c(Q),\qquad
 r_c=\lfloor\alpha L_c(n)\rfloor,\qquad \alpha>0.
\]
Then along any unbounded sequence of such pairs
\begin{equation}\label{eq:packetResponse}
 \frac{q^n L_c(n)^2}{r_c}
       \bigl(R_c^{\rm swapped}-R_c^{\rm original}\bigr)
 \longrightarrow \mathbf1_{[P]=c}-\mathbf1_{[Q]=c}
\end{equation}
in the sense that the difference from the displayed, possibly varying,
right side tends uniformly to zero. Consequently the sum of the
squared normalized responses tends to \(2\).
\item At all full degree cutoffs, with \(L_c=\kappa H_n+C_c\) and
\(r_c=\lfloor\alpha L_c\rfloor\), the same reference obeys
\[
 nq^{n/2}\frac{L_c^2}{r_c}
       (R_c-{\cal R}_{c,r_c}(L_c))
       =-nq^{n/2}E_c(n)+o(1).
\]
Thus the packet retains the critical-scale zero response as well as
the marked replacement response.
\end{enumerate}
\end{theorem}
\begin{proof}
Write
\(L(u,\chi)=\prod_j(1-\beta_{\chi,j}u)\) for \(\chi\ne1\).
For an irreducible modulus these characters are primitive.
The finite-field RH gives \(|\beta_{\chi,j}|\leq\sqrt q\);
even characters retain their root \(1\), rather than silently dropping
it. These standard facts and the explicit formula are recorded in
Keating--Rudnick~\cite{KeatingRudnick}, Section 3.
The principal factor is \((1-u^d)/(1-qu)\).

Define the exact proper-power packet
\[
 V_c(m)=\sum_{\substack{e\mid m\\e<m}}e
       \sum_{\deg P=e}\left(1+
             \mathbf1_{[P^{m/e}]=c}\right).
\]
Orthogonality, with characters extended by zero at \(\mathfrak r\),
and the logarithmic derivative of each \(L\)-factor give
\begin{equation}\label{eq:carrierTrace}
 m\sum_{\deg P=m}(1+\mathbf1_{[P]=c})
 =\kappa q^m-\frac d\phi\mathbf1_{d\mid m}
 -\frac1\phi\sum_{\chi\ne1}\overline{\chi(c)}
           \sum_j\beta_{\chi,j}^{\,m}-V_c(m).
\end{equation}
In particular the proper-power indicator is
\(\mathbf1_{[P^{m/e}]=c}\); replacing it by an indicator of \(P\)
would change the identity.

The degree-\(m\) discrepancy from \(\kappa/m\) is absolutely summable:
the trace term is \(O_d(q^{-m/2}/m)\), the proper-power term has
\(e\leq m/2\), and the overlap correction is \(O(q^{-m}/m)\).
Define \(C_c\) by that convergent full series. Summing its tail gives
the exact identity
\begin{align}
 E_c(n)&=\sum_{m>n}\frac1{m q^m}
 \left\{\frac d\phi\mathbf1_{d\mid m}
 +\frac1\phi\sum_{\chi\ne1}\overline{\chi(c)}
                   \sum_j\beta_{\chi,j}^{\,m}
 +V_c(m)\right\}\nonumber\\
 &\hspace{8mm}+
 \sum_{\substack{\deg P>n\\{[P]=c}}}q^{-2\deg P}.
 \label{eq:carrierClockTail}
\end{align}
For example, its absolute value is at most
\[
 (2d+4)\frac{q^{-(n+1)/2}}{(n+1)(1-q^{-1/2})}
 +\frac{q^{-(n+1)}}{(n+1)(1-q^{-1})}.
\]
Indeed \(\sum_{\deg P=e}e\leq q^e\), and
\(V_c(m)\leq2\sum_{e\leq m/2}q^e\leq4q^{m/2}\).
This proves the uniform clock estimate, retaining every correction.

Subtract the ramification term, \(V_c\) tail, and final overlap tail
in \eqref{eq:carrierClockTail}; call the result \({\cal Z}_c(n)\).
For every \(\chi\ne1\), orthogonality gives exactly
\[
 {\cal F}_\chi(n):=\sum_{c\in G}\chi(c){\cal Z}_c(n)
       =\sum_j\sum_{m>n}\frac{\beta_{\chi,j}^{\,m}}{m q^m}.
\]
Thus
\[
 nq^n\bigl({\cal F}_\chi(n-1)-{\cal F}_\chi(n)\bigr)
       =\sum_j\beta_{\chi,j}^{\,n}.
\]
Newton identities recover the finite \(L\)-polynomial, with multiplicity.
There is no RH hypothesis about the integers in this argument.
The clocks themselves are observables of the full density law:
\(\eta_i=d_1(i)/(1-\sum_{j<i}d_1(j))\).
This uses the existing inverse-recovery mechanism; it is not an
additional zero-recovery assertion about approximate data.

For the marked part, \([P]\ne[Q]\), because \(\mathfrak r\nmid h\).
Their odds in each class are \(O(q^{-n})\). Formula
\eqref{eq:classCarrierProbability} gives
\[
 q^n(w_c(P)-w_c(Q))
 =\mathbf1_{[P]=c}-\mathbf1_{[Q]=c}+O(q^{-n}).
\]
Apply Lemma~\ref{lem:coupledResponse} on the common prefix without
either terminal prime. This proves \eqref{eq:packetResponse} without
subtracting two independent errors on the zero scale.
There are exactly two nonzero coordinates in the limiting class
contrast, equal to \(1,-1\); its squared energy is exactly \(2\).

Finally substitute \(B_c-L_c=E_c(n)\) into
\eqref{eq:markedLinearization}. After the stated normalization its
error is \(O(1/r_c)+O(nq^{-n/2})=o(1)\).
The same calculation applies to the pair cut after subtracting
\(\eta_c(Q)\) from both clocks. The two observations therefore use
one canonical reference, with the geometric labels still attached.
\end{proof}

The theorem separates all existing \(h\)-edges using finitely many
laws with no character blind spot. Its zero readout uses all degree
cutoffs; no assertion is made that an arbitrary sparse subsequence of
edge degrees resolves every Frobenius phase. The vector of residue
contrasts is not, by itself, a unique recovery of the polynomial \(h\);
the exact equation \(Q=P+h\) is retained in the prime labels.

\begin{corollary}[The geometric gap limit from one finite packet]
\label{cor:packetGapLimit}
Fix \(H\geq0\) and choose the modulus above with \(d>H\), \(d\geq2\).
Let \({\cal P}_n\) be the finite set of monic irreducibles of degree \(n\),
\(S_n={\cal P}_1\cup\cdots\cup{\cal P}_n\), and
\[
 {\cal B}_{n,H}=\{(P,Q)\in{\cal P}_n^2:
                   P\ne Q,\ \deg(Q-P)\leq H\}.
\]
For each such pair use
\(L_{c,Q}=\kappa H_n+C_c-\eta_c(Q)\),
\(r_{c,Q}=\lfloor\alpha L_{c,Q}\rfloor\), and define
\[
 U_c(P,Q;n)=\frac{q^n L_{c,Q}^2}{r_{c,Q}}
 \left(R_{c,r_{c,Q}}(S_n\setminus\{P\})
             -R_{c,r_{c,Q}}(S_n\setminus\{Q\})\right).
\]
These quantities are defined for all sufficiently large \(n\).
Set the empty-edge value explicitly by
\[
 {\cal A}_{n,H}=\max\left(\{0\}\cup
       \left\{\sum_{c\in G}U_c(P,Q;n)^2:
                              (P,Q)\in{\cal B}_{n,H}\right\}\right).
\]
If \({\cal B}_{n,H}\) is empty, \({\cal A}_{n,H}=0\).
Otherwise, uniformly in every pair in the finite domain,
\[
 {\cal A}_{n,H}=2+
 O_{q,d,\alpha}\left(\frac1{\log n}+\frac{\log n}{q^n}\right).
\]
Consequently
\[
 \limsup_{n\to\infty}{\cal A}_{n,H}=2
 \quad\Longleftrightarrow\quad
 \text{some fixed nonzero \(h\), \(\deg h\leq H\), recurs infinitely often}.
\]
Equivalently, if
\[
 {\mathfrak H}_q=
 \liminf_{n\to\infty}
     \min_{\substack{P,Q\in{\cal P}_n\\P\ne Q}}\deg(Q-P)
       \ \in\ \mathbb Z_{\geq0}\cup\{\infty\},
\]
where the minimum is needed only at sufficiently large degrees,
then \({\mathfrak H}_q\leq H\) if and only if this limsup is \(2\).
The corresponding least recurring polynomial norm distance is
\(q^{{\mathfrak H}_q}\) when \({\mathfrak H}_q\) is finite.
\end{corollary}
\begin{proof}
A nonzero polynomial of degree at most \(H<d\) cannot be divisible
by \(\mathfrak r\). Thus the same fixed packet separates every pair
in \({\cal B}_{n,H}\), without selecting a different modulus for
each shift. The bounds in Lemma~\ref{lem:coupledResponse} are uniform
over the deleted terminal pair: both odds are \(O(q^{-n})\), the
background clock changes by \(O(q^{-n})\), and all ratio bounds use
the same fixed small-weight tail. Since \(r_{c,Q}\asymp L_{c,Q}
\asymp\log n\), the normalized coordinate error is
\(O(1/\log n+(\log n)/q^n)\). Summing the finitely many squared
coordinates gives the displayed uniform energy estimate and hence
the maximum estimate.

Infinitely many nonempty \({\cal B}_{n,H}\) are equivalent to an
infinite-degree sequence of pairs. There are only
\(q^{H+1}-1\) possible nonzero differences of degree at most \(H\);
one fixed difference therefore occurs on an unbounded sequence.
Conversely, such a recurring difference makes \({\cal B}_{n,H}\)
nonempty infinitely often. The uniform energy estimate proves both
directions of the limsup criterion. The characterization of
\({\mathfrak H}_q\) follows because degrees are nonnegative integers.
\end{proof}

This is a gap-limit criterion using the zero-sensitive packet itself;
it is not a relabelling of the unmarked norm-gap threshold.
For comparison, \({\mathfrak H}_2\geq2\): every irreducible of degree
greater than one takes value \(1\) at both \(0\) and \(1\), so every
difference of two such primes is divisible by \(T(T+1)\).
The next example gives \({\mathfrak H}_7=0\), the least possible value.
The unmarked critical ascent coefficient is nevertheless \(1\) in
both fields. No independence between the different class laws is
used; they can all be sampled from the same pair \((a,b)\).

\subsection{An explicit recurring pair with a surviving elliptic wave}
There is also a fixed polynomial family, with no factorization filter,
in which both channels can be evaluated completely.
Take \(q=7\), \(D=T^3-T\), and sample independent monic \(a,b\) of
degree \(M\) in the norm form
\[
                              F(a,b)=a^2-Db^2.
\]
For a prime of norm \(Q\), put \(\chi=\chi_D(P)\), with
\(\chi=0\) at \(P\mid D\). Direct counting over the residue field gives
\begin{equation}\label{eq:normFormProbability}
                    \eta_F(P)=\frac{1+\chi}{Q}-\frac{\chi}{Q^2}.
\end{equation}
At a split prime the two lines have \(2Q-1\) points in total; at an
inert prime only \((0,0)\) is a zero; at a ramified prime \(a=0\)
and \(b\) is arbitrary. Thus the respective probabilities are
\(2/Q-1/Q^2,\ 1/Q^2,\ 1/Q\), all strictly between zero and one.
CRT again proves the fixed-cutoff independent law before \(M\to\infty\).

The elliptic curve \(y^2=T^3-T\) has trace zero over \(\F_7\):
\(x\) and \(-x\) cancel in the quadratic sum since
\(\chi_7(-1)=-1\). Hence
\[
              L_D(u)=1+7u^2,\qquad \beta_\pm=\pm i\sqrt7.
\]
Let \(C_D\) and \(E_D(n)=B_D(n)-H_n-C_D\) be the quadratic clock
constant and error already defined above. Put
\[
 C_F=C_D-\sum_P\frac{\chi_D(P)}{|P|^2},\qquad
 E_F(n)=B_F(n)-H_n-C_F.
\]
The absolutely convergent overlap correction gives exactly
\[
 E_F(n)=E_D(n)+\sum_{\deg P>n}\frac{\chi_D(P)}{|P|^2}.
\]

\begin{theorem}[Simultaneous gap mark and elliptic zero wave]
\label{thm:explicitMarkedWave}
Put \(n=3^m\), \(Q=7^n\), and
\[
                  A_n=T^n-4,\qquad B_n=T^n-2.
\]
These are irreducible, \(B_n-A_n=2\), and
\(\chi_D(A_n)=1,\ \chi_D(B_n)=-1\).
Place them last in their degree block in this orientation and compare
the intervening prefix with the swapped prefix.
For \(L_n=H_n+C_F-Q^{-2}\),
\(r_n=\lfloor\alpha L_n\rfloor\), and fixed \(\alpha>0\), one has
\begin{align}
 \frac{Q L_n^2}{r_n}(R^{\rm swapped}-R^{\rm original})&\longrightarrow2,
 \label{eq:explicitEdgeLimit}\\
 \frac{n\sqrt Q\,L_n^2}{r_n}
       (R^{\rm original}-{\cal R}_{r_n}(L_n))
       +\frac{\sqrt7}{3}
 &=(-1)^m\frac{\sqrt7}{4}+o(1).
 \label{eq:explicitZeroLimit}
\end{align}
\end{theorem}
\begin{proof}
The binomial irreducibility criterion applies for every \(m\), since
\(2,4\) have order \(3\) in \(\F_7^\times\), and \(n\) is a power
of \(3\). This is the constant-shift construction of
Burrin--Issac~\cite{BurrinIssac}.
Since \(n\) is odd, \(T^n\equiv T\bmod D\).
Quadratic reciprocity, including its odd-degree sign, gives
\(\chi_D(T^n-a)=\chi_7(a^3-a)\).
The values at \(a=4,2\) are respectively
\(\chi_7(4)=1,\chi_7(6)=-1\).

The original prefix clock is \(B_F(n)-Q^{-2}=L_n+E_F(n)\).
The swapped prefix clock is
\(B_F(n)-(2/Q-Q^{-2})=L_n+E_F(n)-2/Q+2/Q^2\).
Their common background is exactly the same. Their inserted odds are
\[
 a=\frac{2Q-1}{(Q-1)^2},\qquad b=\frac1{Q^2-1},\qquad
 a-b=\frac{2Q^2}{(Q-1)^2(Q+1)}.
\]
All denominators are positive, including at the smallest norm \(Q=7\).
Lemma~\ref{lem:coupledResponse} proves \eqref{eq:explicitEdgeLimit}.

To evaluate the other channel, retain the exact prime-power tail.
For odd \(n\), the Frobenius tail is
\(\sum_{j>n/2}(-1)^j/(j7^j)\). The prime-square tail is
\(\sum_{j>n/2}1/(j7^j)+O(7^{-3n/4}/n)\), by the exact prime
polynomial formula. All powers at least three contribute
\(O(7^{-2n/3}/n)\); the finitely many ramified corrections and the
norm-form overlap tail are smaller. Consequently
\[
 E_F(n)=\sum_{j>n/2}\frac{1+(-1)^j}{j7^j}
                     +O(7^{-2n/3}/n).
\]
In particular the cancellations at odd \(j\) have not been discarded.
Write \(j_0=(n+1)/2\). Summing the two geometric tails, using
\(n/(j_0+k)\to2\) with a summable geometric majorant, gives
\[
 n7^{n/2}E_F(n)
   =\frac2{\sqrt7}\left(\frac7{6}
                   +(-1)^{j_0}\frac7{8}\right)+o(1)
   =\frac{\sqrt7}{3}+(-1)^{m+1}\frac{\sqrt7}{4}+o(1).
\]
Now apply \eqref{eq:markedLinearization} with the original clock.
Its normalized error is \(O(1/r_n)+O(n7^{-n/2})\), proving
\eqref{eq:explicitZeroLimit}. The constant \(\sqrt7/3\) is the
retained prime-square bias; the alternating term is the elliptic
Frobenius wave. The change from the two original ratios to their
reference deviation and coupled difference is invertible, so neither
channel is erased.
\end{proof}

\paragraph{What this transfers, and the remaining arithmetic question.}
The canonical reference now coexists with a nonzero geometric shift
response, uniformly in growing rank, in an actual positive law.
The complete finite packet prevents losing an existing edge merely
because one chosen character takes equal values at its endpoints.
The number \(2\) in \eqref{eq:explicitEdgeLimit} is the local odds
contrast, not a deduction that the integer \(H_1\) is \(2\).
For \(\F_7[T]\) this example does prove infinitely many constant
gap-\(2\) pairs, already known arithmetically; their polynomial
norm distance is \(1\), rather than the difference of their norms.
The reference transfer does not assert new priority for this pair
construction or for Fourier inversion.

The full integer ascent equivalence uses the actual recurring-edge
supply as well as its zero-sensitive reference. In the finite-field
model RH proves the clock estimate, and the theorem above preserves
both pieces of information; edge supply is not a consequence of
discarding the proper-power packet. Inferring an integer bound for
\(H_1\) would additionally require an arithmetic comparison of those
joint edge counts across the two settings. No such specialization
or characteristic-zero limit is asserted.

\paragraph{Formalization and verification status.}
The exact scalar and elementary-symmetric coupled responses, strict
positivity, the norm-form odds identity, and the complete packet energy,
signed cross-term expansion and two-coordinate lower bound have
real proofs compiled in
\texttt{Helpers/MarkedCoefficientResponse.lean}, exported by the
maintained project facade. The uniform reference and slope providers
used above were already proved in Lean. The new finite-field sampling
adapters, the analytic reindexing application, the Dirichlet tail
identity, the asymptotic transfer theorems and the geometric gap-limit
corollary are ordinary proofs, not newly compiled Lean theorems. Their concrete
formalization dependencies are polynomial-prime CRT counting,
finite-field Dirichlet explicit formulas, and the polynomial-to-rational
prime index growth comparison. The earlier effective-divisor adapter
and fixed-field critical-scale result also remain ordinary proofs.
No new integer zero-free exponent or bound on the integer \(H_1\)
is claimed.

\section{Coherent accumulation across prime characteristics}
Fix a positive even integer \(h\). We retain this same integer shift
throughout, reducing it modulo each rational prime \(p\).
The integer conclusion sought is infinitely many pairs \(n,n+h\)
of rational primes. It would give \(H_1\leq h\); for \(h=2\) it gives
\(H_1=2\). This is distinct from the degree bound used for
\({\mathfrak H}_q\) above.

\begin{lemma}[Every admissible local residue has a monic lift]
\label{lem:artinSchreierCoverage}
For every prime \(p\) and \(a\in\F_p\) with \(a(a+h)\ne0\),
the monic polynomials
\[
                 T^p-T+a,\qquad T^p-T+a+h
\]
are irreducible over \(\F_p\).
Thus all admissible constant terms can be covered in every prime
characteristic. When \(p\mid h\), the two reductions coincide;
this is the required local overlap, not a distinct-pair assertion.
\end{lemma}
\begin{proof}
If \(\alpha\) is a root of the first polynomial, then
\(\alpha^p=\alpha-a\). Its successive Frobenius images are
\(\alpha-ka\). For \(a\ne0\) their orbit has exactly \(p\) elements,
so the minimal polynomial has degree \(p\).
Apply the same argument to \(a+h\).
\end{proof}

This establishes a genuine compatibility step for a correctly
paritied shift. It does not replace the earlier growing-degree
recurrence and reference information by a parity test.
For a finite prime set, CRT can join whole coefficient lists of
these polynomials, padding with zero coefficients where degrees
differ; it gives a common integer polynomial with the prescribed
reductions. Such a lift need not be monic over the integers when
the local degrees differ. More importantly, CRT alone gives no
small-height bound for the lift or its values.

\subsection{An elliptic curve linking the characteristic-wise counts}
There is a richer local packet than one existence witness.
For an odd prime \(p\), \(h\ne0\) in \(\F_p\), and
\(c(c+h)\ne0\), let \(w_p(c;h)\) count all \(b\in\F_p\) for which
\[
                 T^2+bT+c,\qquad T^2+bT+c+h
\]
are both irreducible. Write \(\chi=\chi_p\), extended by zero,
and let \(a_p(c,h)\) be the Frobenius trace of
\[
                       E_{c,h}:y^2=x(x-c)(x-c-h).
\]
The curve is nonsingular under the stated hypotheses.

\begin{proposition}[The conditional quadratic count]
\label{prop:quadraticEllipticCount}
One has the exact identity
\begin{align}
4w_p(c;h)={}&p+1-a_p(c,h)\nonumber\\
 &-(1+\chi(c))(1-\chi(-h))
  -(1+\chi(c+h))(1-\chi(h)).
 \label{eq:quadraticEllipticCount}
\end{align}
In particular, \(w_p(c;h)>0\) for every admissible \(c\) when
\(p\geq17\).
The total count is
\[
 \sum_{c\in\F_p}w_p(c;h)
       =\frac p4\bigl(p-3+\chi(h)+\chi(-h)\bigr).
\]
Here \(w_p(0;h)=w_p(-h;h)=0\).
\end{proposition}
\begin{proof}
The discriminants are \(b^2-4c\) and \(b^2-4(c+h)\).
First expand the product of their nonsquare indicators with the
zero-discriminant corrections still separate. The two individual
character sums equal \(-1\).
The coupled character sum is
\[
 \sum_b\chi((b^2-4c)(b^2-4(c+h)))
 =\sum_x(1+\chi(x))\chi((x-c)(x-c-h))
 =-1-a_p(c,h).
\]
The substitution uses that \(b^2=4x\) has \(1+\chi(x)\) solutions.
The first term in the last sum is \(-1\), while the second is the
negative elliptic Frobenius trace.

The uncorrected count is therefore
\((p+1-a_p(c,h))/4\).
There are \(1+\chi(c)\) roots with first discriminant zero, and at
each of them the other discriminant has character \(\chi(-h)\).
The other zero-discriminant row similarly contributes
\((1+\chi(c+h))(1-\chi(h))/4\).
The two rows do not intersect because \(h\ne0\).
This proves the full formula, including every overlap correction.

Hasse's bound gives the uniform lower estimate
\[
 w_p(c;h)\geq\frac{p-7-2\sqrt p}{4}>0\qquad(p\geq17).
\]
For the total count, put \(u=b^2-4c\).
For each \(b\), varying \(c\) varies \(u\) bijectively.
Counting \(u,u-4h\) both nonzero nonsquares gives
\((p-3+\chi(h)+\chi(-h))/4\); there are \(p\) choices of \(b\).
\end{proof}

\paragraph{Formal status of this packet.}
The exact count identity, its total over all constants, and the vanishing
at \(c=0,-h\) are proved in maintained Lean over every odd finite field.
The declarations
\texttt{monicQuadraticPairCount\_eq\_trace} and
\texttt{sum\_monicQuadraticPairCount\_eq} in the candidate module
\texttt{LegendreSymbol.QuadraticPolynomialPairs}
use the actual irreducibility predicate and actual cubic point counts.
Their checked axiom closures contain only the foundational Lean axioms.
The Hasse deduction of positivity and the subsequent CRT and
short-interval accumulation assertions still have ordinary proofs here;
they are not claimed as formalized integer prime-pair estimates.

For fixed positive integers \(n,h\), the traces in this proposition
at all primes not dividing \(2hn(n+h)\) are reductions of one
elliptic curve over \(\mathbb Q\):
\[
                         E_{n,h}:y^2=x(x-n)(x-n-h).
\]
Thus these are arithmetically linked field cases.
Their exact count is not just the Boolean statement that a local
pair exists. Elliptic and modular information in finite-field
prime-pair counts also occurs in the cubic computation of
Bary-Soroker--Stix~\cite{BarySorokerStix}; no novelty claim for that
general connection is made here.

\subsection{A concrete integer-side accumulation packet}
For each \(p\geq17\) with \(p\nmid h\), define
\[
 u_p(c)=\frac{p\,w_p(c;h)}{\sum_{a\in\F_p}w_p(a;h)}.
\]
For the finitely many other primes, use the Artin--Schreier packet:
\[
 u_p(c)=\frac{p}{p-\nu_p(h)}\,\mathbf1_{c(c+h)\ne0},
 \qquad
 \nu_p(h)=\begin{cases}1,&p\mid h,\\2,&p\nmid h.\end{cases}
\]
Evenness of \(h\) ensures \(p-\nu_p(h)>0\) also at \(p=2\).
For every prime,
\[
 \frac1p\sum_{c\in\F_p}u_p(c)=1,\qquad
 u_p(c)>0\ \Longleftrightarrow\ c(c+h)\ne0.
\]
The positivity assertion uses the preceding proposition at large
primes and the explicit construction at every exceptional prime.

For good \(p\) and admissible \(c\), the exact signed form is
\[
 u_p(c)=1+
 \frac{2-a_p(c,h)-\chi(c)(1-\chi(-h))
                   -\chi(c+h)(1-\chi(h))}
      {p-3+\chi(h)+\chi(-h)}.
\]
Equivalently, writing
\(Z_p(c,h)=(1+\chi(c))(1-\chi(-h))
 +(1+\chi(c+h))(1-\chi(h))\), it is the finite local quotient
\[
 \frac{1-a_p(c,h)/p+1/p}{1-(3-\chi(h)-\chi(-h))/p}
        \left(1-\frac{Z_p(c,h)}{p+1-a_p(c,h)}\right).
\]
The first numerator is the inverse good-prime Euler factor of
\(L(E_{n,h},s)\) at \(s=1\), when \(c=n\bmod p\).
All products below are finite. No central Euler-product limit or
zero-recovery theorem for elliptic \(L\)-functions is assumed.
Hasse's bound is used to certify local positivity; the proposed
global estimate retains the signed traces and compensation factors.

\begin{proposition}[Exact accumulation and its required height scale]
\label{prop:integerAccumulation}
For real \(Y\geq2\), set
\[
 M_Y=\prod_{p\leq Y}p,\qquad
 U_Y(n)=\prod_{p\leq Y}u_p(n\bmod p).
\]
Then
\[
 \sum_{n=1}^{M_Y}U_Y(n)=M_Y,\qquad
 U_Y(n)>0\ \Longleftrightarrow\
                \gcd(n(n+h),M_Y)=1.
\]
For an integer \(X\geq h+3\), put \(Y=\sqrt{2X+h}\).
There is a prime pair \(n,n+h\) with \(X<n\leq2X\) if and only if
\[
                         S_h(X):=\sum_{X<n\leq2X}U_Y(n)>0.
\]
Consequently integer recurrence follows if this positivity is
established along any unbounded sequence of \(X\).
\end{proposition}
\begin{proof}
CRT, not an independence assumption, gives
\[
 \sum_{n\bmod M_Y}\prod_{p\leq Y}u_p(n\bmod p)
       =\prod_{p\leq Y}\sum_{c\bmod p}u_p(c)
       =\prod_{p\leq Y}p=M_Y.
\]
The support statement follows from the exact local support in
each factor. In the specified interval, \(n,n+h>Y\) and both
are at most \(Y^2\). A composite value would have a prime divisor
at most its square root, hence at most \(Y\).
Thus positive weight forces both values prime; conversely each
prime pair has positive weight. Nonnegativity proves the equivalence.
\end{proof}

This gives a precise meaning to accumulating enough field cases:
they must control the joint integer mass at \(X\) of order \(Y^2\),
rather than only over the complete period
\(M_Y=\exp((1+o(1))Y)\).
The full-period identity is proved; no critical short-interval
lower bound is being inferred from it.

To retain every phase, define
\[
 \widehat u_p(k)=\frac1p\sum_{a\bmod p}u_p(a)e^{-2\pi i ka/p},
 \qquad D_X(t)=\sum_{X<n\leq2X}e^{2\pi i nt}.
\]
Fourier expansion gives the exact finite signed consumer
\[
 S_h(X)=X+
 \sum_{\substack{k\in\prod_{p\leq Y}\F_p\\
                  k\ne0}}
 \left(\prod_{p\leq Y}\widehat u_p(k_p)\right)
 D_X\left(\sum_{p\leq Y}\frac{k_p}{p}\right).
\]
Representatives \(0,\ldots,p-1\) may be used for each \(k_p\).
The sum is real by conjugate pairing. Positivity requires only that
its retained nonprincipal contribution exceed \(-X\), along
unbounded cutoffs; neither a triangle bound on its terms nor a
prime-pair asymptotic is assumed. The characteristic-wise traces
come from the same global curve for a fixed integer \(n\), so the
additional arithmetic information is available in this packet.
Estimating this signed incomplete packet remains open here.
These identities do not establish a new integer gap bound.

\paragraph{Verification status for this section.}
The exact quadratic count, total count, exceptional constants,
normalization, and its rational Euler factorization have maintained
Lean proofs.
The elliptic comparison and its finite-product estimates are formalized
with their explicit numerical trace bounds. The exact moment, centered
variance and finite product-space energy theorems below are formalized
without any trace-bound hypothesis.
The classical Hasse input, Artin--Schreier coverage, finite Fourier
and CRT assertions, and integer short-interval consumer still have
ordinary proofs here. These are separate statuses.

\subsection{A uniform elliptic--Dirichlet factorization}
Fix positive integers \(n,h\), with \(h\) even. All products indexed
by \(p\) in this subsection are over rational primes. At a prime
\(p\geq17\) not dividing \(hn(n+h)\), put
\[
 a=a_p(E_{n,h}),\qquad
 A_p=1-\frac ap+\frac1p,
\]
and, with \(\chi=\chi_p\), define
\[
 Q_p=\frac{(1-\chi(n)/p)(1-\chi(n+h)/p)}
 {(1-1/p)(1-\chi(-hn)/p)(1-\chi(h(n+h))/p)}.
\]
For each of the seven integer labels
\(1,n,n+h,h,-h,-hn,h(n+h)\), write
\[
 B_{\xi,p}=1-\frac{\chi(\xi)a}{p^2}+\frac1{p^3},\qquad
 T_p=\frac{B_{1,p}^{\,2}B_{n,p}B_{n+h,p}}
 {B_{h,p}B_{-h,p}B_{-hn,p}B_{h(n+h),p}}.
\]
Labels representing the same quadratic character retain their
displayed multiplicities. The factor \(A_p\) is the inverse elliptic
Euler factor at \(s=1\); \(B_{\xi,p}\) is the inverse factor of its
quadratic twist at \(s=2\). The factor \(Q_p\) retains the displayed
quadratic Dirichlet and zeta factors. All formulas here use finite
Euler products, including at \(s=1\).

\begin{theorem}[Uniform seven-twist comparison]
\label{thm:sevenTwistComparison}
At every such prime there is an explicitly defined \(G_p(n,h)>0\)
such that
\[
 u_p(n\bmod p)=A_pQ_pT_pG_p(n,h),\qquad
 |G_p(n,h)-1|\leq\frac{32832}{p^2}.
\]
For all integer cutoffs \(256\leq Y\leq Z\),
\[
 \exp\left(-\frac{65664}{Y}\right)
 \leq
 \prod_{\substack{Y<p\leq Z\\p\nmid hn(n+h)}}G_p(n,h)
 \leq
 \exp\left(\frac{65664}{Y}\right).
\]
The constants are independent of \(n,h,Y,Z\).
\end{theorem}
\begin{proof}
Write
\[
 e=\chi(h),\quad f=\chi(-h),\quad x=\chi(n),\quad y=\chi(n+h),
 \quad d=3-e-f,
\]
and retain the entire zero-discriminant correction
\[
 Z_p=(1+x)(1-f)+(1+y)(1-e).
\]
Each sign is \(1\) or \(-1\), and \(Z_p\in\{0,4,8\}\).
The exact count and total count give
\[
 u_p(n)=\frac{p+1-a-Z_p}{p-d}.
\]
Hasse gives \(a^2\leq4p\), hence \(|a|\leq p/2\) for \(p\geq16\).
Define
\begin{align*}
 N&=(p+1-a-Z_p)(p-1)(p-fx)(p-ey),\\
 D&=(p-d)(p+1-a)(p-x)(p-y),\qquad R=N/D.
\end{align*}
Then \(u_p=A_pQ_pR\), \(N,D>0\), and \(D\geq p^4/16\).

Here is an explicit bound on the signed defect. Put
\[
 U=fx+ey,\quad V=efxy,\quad
 H=U+V-d(x+y)-xy,\quad J=xy(d-ef).
\]
Expansion, before estimating any term, gives
\begin{align*}
 N-D+Z_pa p^2
  ={}&[Z_p(2+U)+H]p^2\\
    &+[-Z_p(U+V)+(1-a)H+J]p+Z_pV+(1-a)J.
\end{align*}
Using \(|Z_p|\leq8\), \(|U|\leq2\), \(|V|\leq1\),
\(|H|\leq14\), and \(|J|\leq6\), its absolute value is at most
\[
 46p^2+(44+14|a|)p+14+6|a|
 \leq53p^2+47p+14\leq56p^2.
\]
The exact four-factor subtraction for \(D-p^4\) gives
\[
 |D-p^4|\leq(47+4|a|)p^3.
\]
Consequently, retaining the compensation cross term,
\begin{align*}
 |(N-D)p^2+Z_paD|
 &\leq56p^4+8|a|(47+4|a|)p^3\\
 &\leq372p^4\leq1024p^4.
\end{align*}
Division by \(Dp^2\), which is positive, proves
\[
 \left|R-1+\frac{Z_pa}{p^2}\right|\leq\frac{16384}{p^2}.
\]

For the shifted factors put \(r=a/p^2\), \(s=p^{-3}\), and
\(k=Z_p/2\in\{0,2,4\}\). Character multiplication gives the exact identity
\[
 T_p=\left(\frac{1-r+s}{1+r+s}\right)^k.
\]
In particular all seven labelled factors, with their signs and
multiplicities, occur in this simplification.
If \(v=(1-r+s)/(1+r+s)-1\), then
\[
 |v|\leq3|r|\leq\frac{3}{2p}\leq\frac18,\qquad
 |v+2r|\leq\frac3{p^2}.
\]
For \(k=0,2,4\), expansion gives
\[
 |(1+v)^k-1-kv|\leq8v^2.
\]
It follows that
\[
 \left|T_p-1+\frac{Z_pa}{p^2}\right|\leq\frac{32}{p^2},
 \qquad T_p\geq(7/8)^4>\frac12.
\]
Thus \(G_p=R/T_p>0\), and the two retained trace terms cancel:
\[
 |G_p-1|\leq
 2\left(\frac{16384}{p^2}+\frac{32}{p^2}\right)
 =\frac{32832}{p^2}.
\]
Only these already compensated remainders were estimated separately.

For \(p>Y\geq256\), the error is at most \(1/2\).
For every \(t\geq1/2\), \(|\log t|\leq2|t-1|\), so
\[
 \left|\log\prod_{\substack{Y<p\leq Z\\p\nmid hn(n+h)}}G_p\right|
 \leq65664\sum_{m=Y+1}^{Z}\frac1{m^2}
 \leq\frac{65664}{Y}.
\]
The last inequality follows by telescoping
\(m^{-2}\leq(m-1)^{-1}-m^{-1}\).
Exponentiating proves the stated finite-tail comparison.
\end{proof}

The general Euler-factor extraction method is classical;
see Alberts~\cite{Alberts}. Here the specified conditional-count packet,
both exceptional discriminant rows, all seven twist labels, and constants
uniform in the integer lift are retained in one comparison.

The complete integer packet still contains its exact opening mask:
\[
 \Gamma_{Y,h}(n)=\mathbf1_{\gcd(n(n+h),M_Y)=1}.
\]
With
\[
 B_h(Y)=\prod_{\substack{p\leq Y\\p<17\ {\rm or}\ p\mid h}}
             \frac p{p-\nu_p(h)},\qquad
 {\cal P}_Y(n,h)=\{p\text{ prime}:17\leq p\leq Y,\ p\nmid hn(n+h)\},
\]
the full product identity is
\[
 U_Y(n)=\Gamma_{Y,h}(n)B_h(Y)
           \prod_{p\in{\cal P}_Y(n,h)}A_pQ_pT_pG_p(n,h).
\]
At a prime dividing \(n(n+h)\), its actual local count is zero.
At a small prime or a prime dividing \(h\), the previously defined
admissible fallback is used. This preserves every support restriction
in the signed short-interval consumer. In particular, the new
comparison supplies a uniform finite-cutoff reference for that consumer;
its short-interval mass remains the arithmetic quantity to estimate.

\paragraph{Formal status of the comparison.}
The maintained modules \texttt{QuadraticPairEuler},
\texttt{QuadraticPairEulerCount}, \texttt{QuadraticPairTwists}, and
\texttt{QuadraticPairProduct} prove the scalar bounds, the exact
seven-factor identity, the actual count factorization, and the finite-tail
comparison. Their arithmetic application uses the explicitly stated
numerical hypothesis \(a^2\leq4p\). The classical Hasse theorem supplies
it in the ordinary finite-field argument above; a Lean proof of Hasse's
theorem is not supplied by these modules. The new declaration closures
use only Lean's foundational axioms.


\subsection{Exact packet energy and its cubic Mertens reference}

This calculation concerns the actual conditional quadratic-pair weights,
rather than an independently chosen random model. Let \(F\) be a finite
field of odd cardinality \(q\), let \(h\in F\setminus\{0\}\), and put
\[
 e=\chi(h),\quad f=\chi(-h),\quad d=3-e-f\in\{1,3,5\}.
\]
Write \(w(c)=w_q(c;h)\) for the irreducible-pair count above.

\begin{theorem}[Exact normalized energy]\label{thm:quadraticPairEnergy}
For every such field and shift,
\[
 16\sum_{c\in F}w(c)^2
   =(q-d)\bigl(q^2+(3-d)q+d\bigr).
\]
If \(q=d\), every count \(w(c)\) is zero. If \(q\ne d\), define
\(u(c)=qw(c)/\sum_{t\in F}w(t)\). Then
\[
 \frac1q\sum_{c\in F}u(c)=1,\qquad
 \frac1q\sum_{c\in F}u(c)^2
   =1+\frac{3q+d}{q(q-d)},
\]
and
\[
 \frac1q\sum_{c\in F}(u(c)-1)^2
    =\frac{3q+d}{q(q-d)}
    =\frac3q+\frac{4d}{q(q-d)}.
\]
In particular \(q\geq16\) excludes a zero normalizer. No Hasse
bound or hypothesis about zeta zeros is required.
\end{theorem}

\begin{proof}
Put \(a_c=a(c,h)\) and \(\kappa_h(x)=\chi(x(x-h))\).
Translation in the defining cubic character sum gives the exact convolution
\[
 a_c=-\sum_{x\in F}\kappa_h(x)\chi(c+x).
\]
The complete two-point character correlation is
\[
 \sum_{c\in F}\chi(c+x)\chi(c+y)=q\mathbf1_{x=y}-1.
\]
Consequently, for every integer-valued function \(b\) on \(F\),
\[
 \sum_{c\in F}\left(\sum_{x\in F}b(x)\chi(c+x)\right)^2
    =q\sum_{x\in F}b(x)^2-\left(\sum_{x\in F}b(x)\right)^2.
\]
Here \(\sum_x\kappa_h(x)=-1\) and
\(\sum_x\kappa_h(x)^2=q-2\). Retaining the mixed correlations as
well as the diagonal therefore yields
\[
 \sum_ca_c=0,\qquad \sum_ca_c^2=q^2-2q-1,\qquad
 \sum_ca_c\chi(c)=\sum_ca_c\chi(c+h)=-1.
\]
The two singular traces are \(a_0=f\) and \(a_{-h}=e\).

Let \(A=q-1+e+f=q+2-d\), \(\alpha=1-f\), \(\beta=1-e\), and
\[
 V_c=A-a_c-\alpha\chi(c)-\beta\chi(c+h).
\]
The exact count formula gives \(4w(c)=V_c\) outside \(c=0,-h\).
At both excluded constants \(w(c)=0\) and \(V_c=q\). Thus the
pointwise square completion is
\[
 16w(c)^2=V_c^2-q^2\mathbf1_{c=0}-q^2\mathbf1_{c=-h}.
\]
Using also \(\sum_c\chi(c)^2=\sum_c\chi(c+h)^2=q-1\) and
\(\sum_c\chi(c)\chi(c+h)=-1\), its sum is exactly
\[
\begin{split}
 16\sum_cw(c)^2={}&qA^2+(q^2-2q-1)
       +(\alpha^2+\beta^2)(q-1)\\
       &-2\alpha-2\beta-2\alpha\beta-2q^2\\
   ={}&(q-d)\bigl(q^2+(3-d)q+d\bigr).
\end{split}
\]
The last equality uses \(e^2=f^2=1\). In particular neither mixed
trace-character term, the character-character term, nor either
exceptional square has been discarded.

When \(q=d\), the sum of the nonnegative squares \(w(c)^2\) is zero,
so all counts vanish. Otherwise the already proved total
\(4\sum_cw(c)=q(q-d)\) gives \(u(c)=4w(c)/(q-d)\).
Substitution proves the second moment and mean. Expanding
\((u(c)-1)^2\) gives the centered formula.
\end{proof}

For real \(q\geq16\) and \(1\leq d\leq5\), define
\[
 m(q,d)=1+\frac{3q+d}{q(q-d)},\qquad
 B(q,d)=m(q,d)(1-1/q)^3.
\]
\begin{theorem}[Uniform finite-packet energy comparison]
\label{thm:quadraticPacketEnergy}
Let \(F_i\), \(i\in I\), be a finite family of odd finite fields,
with \(q_i\geq16\), and let \(h_i\ne0\).
On \(\Omega=\prod_{i\in I}F_i\) with its uniform measure, put
\(W(x)=\prod_i u_{h_i}(x_i)\) and \(d_i=3-\chi_i(h_i)-\chi_i(-h_i)\).
Then
\[
 \mathbb E_\Omega W^2=\prod_{i\in I}m(q_i,d_i),\qquad
 \mathbb E_\Omega W^2\prod_{i\in I}(1-1/q_i)^3
       =\prod_{i\in I}B(q_i,d_i).
\]
Every correction is positive and satisfies
\[
                 |B(q,d)-1|\leq\frac{20}{q^2}.
\]
More explicitly, let \(16\leq Y\leq Z\) be integers and let
\(T\subseteq\{p\text{ prime}:Y<p\leq Z\}\).
Choose any \(d_p\in[1,5]\). Then
\[
 e^{-40/Y}\leq\prod_{p\in T}B(p,d_p)\leq e^{40/Y}.
\]
This includes the actual conditional-count factors with a fixed even
integer \(h\), after restricting \(T\) to \(p\nmid h\).
\end{theorem}

\begin{proof}
Finite product summation, applied to \(u_{h_i}^2\), factors the
expectation exactly. The first theorem gives each factor \(m(q_i,d_i)\);
multiplication by the stated cubic factors gives the second identity.

For completeness, all constants in the remainder estimate can be checked
by signed polynomial identities. Put
\[
 R=(4d-6)q^3+(8-6d)q^2+(4d-3)q-d.
\]
Direct algebra gives \(B-1=R/[q^4(q-d)]\). The two bounds follow from
\[
 R+2q^2(q-d)
   =4(d-1)q^2(q-2)+(d-1)(4q-1)+(q-1)\geq0
\]
and
\[
\begin{split}
 20q^2(q-d)-R={}&6q^2(q-16)+18q(q-1)+q+5\\
       &+(5-d)(4q^3+14q^2+4q-1)\geq0.
\end{split}
\]
Since \(q^4(q-d)>0\), these imply \(|B-1|\leq20/q^2\).
Positivity follows directly from the definition of \(B\).

For \(p>Y\geq16\), this error is less than \(1/2\).
Hence \(|\log B(p,d_p)|\leq40/p^2\). Summing over the actual
finite support \(T\) and using
\(\sum_{Y<n\leq Z}n^{-2}\leq1/Y\) gives the asserted
logarithmic bound and then the exponential inequalities.
\end{proof}

\paragraph{What accumulation adds.}
For the eligible primes the exact local energy is
\[
 m(p,d_p)=1+\frac3p+\frac{4d_p}{p(p-d_p)}.
\]
Thus its complete finite-product reference is the cube of the Mertens
product, with a positive correction having the explicit tails just proved.
For comparison, the normalized flat weight on the \(p-2\) admissible
constants has second moment \(p/(p-2)=1+2/p+O(p^{-2})\).
The conditional irreducibility weights have one additional harmonic
unit of full-period energy. This quantifies the extra elliptic
fluctuation; it is not an impossibility result for signed short-interval
cancellation. Integer recurrence still requires the specific mass
in the earlier window \(X<n\leq2X\) at cutoff \(Y\asymp\sqrt X\).
A product-space or full-period energy calculation alone does not supply
that window estimate.

\paragraph{Scope and formal status.}
Moments of the Legendre family are classical; see, for example,
Kane--Zindulka~\cite{KaneZindulka}. The trace second moment itself is
not asserted as new. The result here computes the actual normalized
conditional-pair statistic, retains its two singular corrections and
all mixed terms, and supplies a uniform finite-packet comparison.
External priority for that combination has not been established.

The maintained module \texttt{QuadraticPairMoments} proves the trace
convolution and moments, the completed count-square identity, the
normalized mean and variance, and the zero-normalizer alternative.
The module \texttt{QuadraticPairEnergy} proves finite product-space
energy factorization, its actual-field specialization, the scalar
\(20/q^2\) bound, and the \(e^{\pm40/Y}\) finite-tail comparison.
These are real Lean proofs without a Hasse hypothesis.
The existing CRT interpretation remains as previously stated; no
new integer prime-pair lower bound is claimed.

\section{Unconditional transfer to quadratic number fields}

There is a direct arithmetic counterpart of finite-word realization.
Let \(\Delta\) be a fundamental quadratic discriminant. Enumerate in
increasing order the odd, unramified rational primes
\(p_1<p_2<\cdots\) splitting in \(\mathbf Q(\sqrt\Delta)\).
For \(c=1\), let \(b_{\Delta,k}^{(1)}(i)\) be the natural density of
integers \(n\) whose \(k\)-th distinct prime factor on this split list
is \(p_i\). For \(c=2\), let \(b_{\Delta,k}^{(2)}(i)\) be the density
of positive integers \(a\) for which \(p_i\) is the \(k\)-th distinct
factor on this list of \(|a^2-\Delta|\). Factors outside the list
do not enter either rank. The prime \(2\) and all ramified primes
are excluded, including from the \(c=2\) root-count formula.

CRT gives the local probabilities \(c/p_i\), and hence the exact formula
\[
 b_{\Delta,k}^{(c)}(i)=\frac c{p_i}
       \prod_{j<i}\left(1-\frac c{p_j}\right)
       e_{k-1}\left(\frac c{p_1-c},\ldots,\frac c{p_{i-1}-c}\right).
                                                               \tag{26}
\]
For \(c=2\) the two residue roots exist precisely because \(p_i\)
splits and is odd and unramified. All local probabilities are less
than one. For \(k\geq2\) and \(i\geq k\), writing \(r=k-1\) and
\(R_r=e_{r-1}/e_r\) on the prefix before \(p_i\), the exact step is
\[
 \frac{b_{\Delta,k}^{(c)}(i+1)}{b_{\Delta,k}^{(c)}(i)}
       =\frac{p_i-c+cR_r}{p_{i+1}},\qquad
 \text{strict ascent}\ \Longleftrightarrow\
       R_r>\frac{p_{i+1}-p_i+c}{c}.                 \tag{27}
\]

\begin{theorem}[Simultaneous number-field word realization]
\label{thm:numberFieldWords}
Fix \(k\geq2\) and a finite word in \(\{+,-\}\).
There are infinitely many real quadratic fields and infinitely many
imaginary quadratic fields such that both integer profiles
\(b_{\Delta,k}^{(1)}\) and \(b_{\Delta,k}^{(2)}\) realize that word
on the same contiguous prime-index segment. This is unconditional.

For each fixed quadratic field both profiles eventually strictly
decrease. Thus their number of positive peak plateaus is finite for
each field and unbounded across either the real or imaginary families.
For \(k=1\) both profiles are strictly decreasing in every field.
\end{theorem}
\begin{proof}
Let the word have length \(t\), put \(r=k-1\), \(s=2r+1+t\),
and \(\varepsilon=r/(8s)\). For all sufficiently large \(Y\geq32s\),
every interval
\[
 J_j(Y)=[4^jY,\,4^jY+\varepsilon Y],
                   \qquad 0\leq j\leq t,
\]
contains at least \(s\) primes, by the ordinary prime number theorem
with \(r,t\) fixed. Choose \(2r+1\) initial primes from \(J_0(Y)\).
For each successive plus select another prime in the current interval;
for each minus select one in the next interval. There are \(s\)
selected primes in total; enough are available in each interval.

For both \(c=1,2\), every selected prefix has total odds at most
\(cs/(Y-c)\), and therefore
\[
 R_r\geq\frac{r(Y-c)}{cs}.
\]
Within a plus step the gap is at most \(\varepsilon Y\).
The choice \(Y\geq32s\) ensures
\[
 \frac{r(Y-c)}s>\varepsilon Y+c \qquad(c=1,2),
\]
so (27) gives a strict ascent in both profiles. At a minus step,
the next prime \(q\) is at least twice the current prime \(p\).
There are at least \(2r\) preceding selected primes. Since their
smallest odds are at least \(c/(p-c)\), (14) gives
\[
 R_r\leq\frac{r(p-c)}{c(m-r+1)}
       \leq\frac{r(p-c)}{c(r+1)}
       <\frac{p-c}{c}<\frac{q-p+c}{c}.
\]
Thus both profiles strictly descend. The initial \(2r+1\) entries
put the entire specified segment inside the positive support.

It remains to realize exactly this list of split primes through its
last element \(B\). For every odd prime \(p\leq B\), prescribe
\(\sigma_p=+1\) on the selected list and \(-1\) elsewhere.
CRT gives a reduced residue class of integers \(\ell\), modulo
\(8\prod_{\substack{p\leq B\\p\ {\rm odd\ prime}}}p\), such that
\[
 \ell\equiv5\pmod8,\qquad (\ell/p)=\sigma_p .
\]
Dirichlet's theorem supplies infinitely many primes \(\ell>B\)
in this class~\cite{Sutherland}. Then \(\Delta=\ell\) is a positive
fundamental discriminant with the prescribed splitting pattern.

For imaginary fields instead prescribe
\[
 \ell\equiv3\pmod8,\qquad
       (\ell/p)=\sigma_p(-1/p).
\]
Again Dirichlet gives infinitely many primes in this reduced class,
and \(\Delta=-\ell\) is a negative fundamental discriminant with
\((\Delta/p)=\sigma_p\). In both constructions \(\Delta\equiv5\pmod8\);
the prime \(2\) remains excluded from the sampling rank. The finite
density formula (26) depends only on the prescribed initial split list,
so every resulting field realizes the same word in both profiles.

Finally the unconditional Mertens theorem for Chebotarev prime sets
gives, for each fixed \(\Delta\),
\[
 \mu_c(x):=
 \sum_{\substack{p\leq x,\ p\ {\rm odd}\\\chi_\Delta(p)=1}}
          \frac c{p-c}
       =\frac c2\log\log x+C_{\Delta,c}+o_\Delta(1).
                                                               \tag{28}
\]
Indeed the split class has density \(1/2\), and replacing \(c/p\)
by \(c/(p-c)\) changes the sum by an absolutely convergent
correction~\cite{AKK}. Hence \(\mu_c\to\infty\).
Since every individual odds weight is at most two, (25) gives
\(R_r\leq r/(\mu_c-2(r-1))\to0\) on sufficiently long prefixes.
Equation (27) forces strict decrease thereafter.
Alternating words \(+,-,+,-,\ldots\) give arbitrarily many strict
peaks across either family. For \(k=1\), the exact ratio is
\((p_i-c)/p_{i+1}<1\).
\end{proof}

\paragraph{Relation to zeta functions and the finite-field result.}
For \(K=\mathbf Q(\sqrt\Delta)\), the standard factorization
\(\zeta_K(s)=\zeta(s)L(s,\chi_\Delta)\) identifies the quadratic
character controlling these splitting profiles with the quadratic
field's zeta factor. The finite-word theorem itself uses only CRT,
the prime number theorem, Dirichlet's theorem, and the coefficient
comparisons; it assumes neither RH nor GRH.

The sharp function-field peak bound above is expressed in the
conductor degree \(\deg D\). That corresponds to a logarithmic
conductor scale, since its finite norm is \(q^{\deg D}\).
An analogous sharp discriminant-uniform bound over \(\mathbf Q\)
has not been proved here. The integer step (27) retains the actual
rational-prime gaps, whereas a fixed polynomial degree block has
equal norms. The fixed-field asymptotic (28) has a field-dependent
constant and remainder; it is not a uniform square-root estimate.
These transfer statements have ordinary proofs, not new Lean
formalizations; no new bound on number-field zeta zeros is claimed.

\section{Arithmetic transfers and the next ascent investigation}
\paragraph{Existing ingredients.}
The generic greedy record sandwich and the maintained periodic and
joint-wave results motivate this argument. The new arithmetic input
in this setting is the established finite-field RH, which supplies
actual bounded Frobenius phases. The prime-power inversion identifies
the reference correction before the signed record consumer is used.
Integer-prime sampling and probability laws still require explicit
closed-point or polynomial adapters.

\paragraph{Why positive genus matters.}
The rational function field has trivial zeta numerator. Nontrivial
zero signals require a positive-genus curve, or a nontrivial polynomial
Dirichlet \(L\)-function. Merely replacing integers by \(\F_q[T]\)
would discard the nontrivial spectral information.

\paragraph{Character and Carmichael structure.}
For squarefree \(f=\prod P_i\in\F_q[T]\), CRT gives
\[
 (\F_q[T]/(f))^\times\simeq
 \prod_i\F_{q^{\deg P_i}}^\times,\qquad
 \operatorname{exponent}((\F_q[T]/(f))^\times)
       =\operatorname{lcm}_i(q^{\deg P_i}-1).
\]
This is the correct counterpart of the integer Carmichael exponent.
Polynomial Carmichael objects already have a literature~\cite{BaeHuSha};
their existence and their Korselt criterion are not new here.

\paragraph{Ascent and constellation status.}
The exact ordinary-density ascent identity, complete unimodality
classification, fixed-rank large-field mode count, quadratic finite-word
realization, and sharp logarithmic conductor-degree mode complexity
have ordinary proofs above. A general marked RH ascent equivalence
and a prime-constellation recurrence transfer are not proved here.
Fixed \(q\) with growing degree and growing \(q\) at fixed degree remain
different limits. The exact quadratic-pair count and total are now
formalized as described above. The full Frobenius sampling and
probability-law adapters remain to be implemented.

\paragraph{Novelty boundary.}
Frobenius-wave descriptions of function-field prime races are established
in the literature~\cite{CFJ,BDKL}. The object here is the minimum irreversible
deletion enforcing every past one-sided cutoff, with its full backward
envelope and elliptic classification. No matching theorem has yet been
identified, but this is not an established priority claim.

\section{Exact diagnostic checks}
The companion script \texttt{finite\_field\_greedy\_probe.py} checks all
general Weierstrass equations over \(\F_2,\F_3\), all short Weierstrass
equations over \(\F_5,\F_7,\F_{11}\), and base extensions of degrees
two and four of the first two families. Singular curves are excluded
by their discriminant. Extension traces satisfy
\(s_0=2,\ s_1=t,\ s_m=ts_{m-1}-qs_{m-2}\).

The domain contains 350 nonsingular equations and 59 realized
\((q,t)\) pairs. Every degree \(1\leq n\leq192\) is checked for
closed-point integrality and nonnegativity, the full identity (3),
and the greedy record sandwich (1).
\begin{center}
\begin{tabular}{c|r|r|r|r}
\(q\)&\(t\)&last deletion through degree 192&deleted mass&deleted count\\ \hline
2&-2&1&3&3\\
2&-1&1&2&2\\
2&0&2&3&2\\
3&0&2&5&3\\
16&8&none&0&0\\
81&18&none&0&0
\end{tabular}
\end{center}
In the other 53 tested pairs, the last observed deletion is at degree
at least 183. These observations neither prove infinite deletion nor
certify a final deletion degree. Universal conclusions above come from
the ordinary proofs, not this finite computation.

\begin{thebibliography}{99}
\bibitem{Milne} J. S. Milne, \emph{Elliptic Curves}, second edition.
\url{https://www.jmilne.org/math/Books/EC2.pdf}.
\bibitem{CFJ} B. Cha, D. Fiorilli, F. Jouve,
\emph{Prime number races for elliptic curves over function fields},
arXiv:1502.05295. \url{https://arxiv.org/abs/1502.05295}.
\bibitem{BDKL} A. Bailleul, L. Devin, D. Keliher, W. Li,
\emph{Exceptional biases in counting primes over function fields},
arXiv:2302.13665. \url{https://arxiv.org/abs/2302.13665}.
\bibitem{BaeHuSha} S. Bae, S. Hu, M. Sha,
\emph{On the Carmichael rings, Carmichael ideals and Carmichael polynomials},
arXiv:1809.05432. \url{https://arxiv.org/abs/1809.05432}.
\bibitem{Cambie} S. Cambie,
\emph{Resolution of Erdos' problems about unimodularity},
arXiv:2501.10333. \url{https://arxiv.org/abs/2501.10333}.
\bibitem{Jung} H. Jung,
\emph{Mertens products in arithmetic progressions over function fields},
arXiv:2602.05788v2. \url{https://arxiv.org/abs/2602.05788}.
\bibitem{Sutherland} A. V. Sutherland,
\emph{Number Theory I}, MIT 18.785, Fall 2021 lecture notes.
\bibitem{AKK} S. Arango-Pineros, D. Keliher, C. Keyes,
\emph{Mertens' theorem for Chebotarev sets},
arXiv:2103.14747. \url{https://arxiv.org/abs/2103.14747}.
\bibitem{BurrinIssac} C. Burrin, M. Issac,
\emph{Infinitely many twin prime polynomials of odd degree},
arXiv:2101.08038. \url{https://arxiv.org/abs/2101.08038}.
\bibitem{KeatingRudnick} J. P. Keating, Z. Rudnick,
\emph{The variance of the number of prime polynomials in short intervals
and in residue classes}, arXiv:1204.0708.
\url{https://arxiv.org/abs/1204.0708}.
\bibitem{KaneZindulka} B. Kane, M. Zindulka,
\emph{Negative bias in moments of the Legendre family of elliptic curves},
arXiv:2608.00435.
\url{https://arxiv.org/abs/2608.00435}.
\bibitem{Alberts} B. Alberts,
\emph{Explicit Analytic Continuation of Euler Products},
Essential Number Theory 5 (2026), 49--112.
\url{https://arxiv.org/abs/2406.18190}.
\bibitem{BarySorokerStix} L. Bary-Soroker, J. Stix,
\emph{Cubic twin prime polynomials are counted by a modular form},
arXiv:1711.05564. \url{https://arxiv.org/abs/1711.05564}.
\end{thebibliography}
\end{document}
